
\newpage
\section{\texorpdfstring{The pre-analytic rings $A_\blacksquare$ and $(A, R)_\blacksquare$}{Solid Rings}}
\label{section: the pre-analytic rings Ablack and (A,R)black}

In this section we will introduce the pre-analytic rings $A_\blacksquare$ and $(A, R)_\blacksquare$ and work towards proving their analyticity. On the way we will try to develop an intuition for $\Z_\blacksquare$-completeness. Before defining these pre-analytic rings we however need to make a technical remark that Clausen and Scholze do not address in \cite{condenseddotpdf}.
\begin{remark}[Natural decomposition of a profinite set]
	To define the various free completions we will break down an extremally disconnected set (which we recall is profinite) into finite constituents $S = \limit_i S_i$. To ensure that the resulting completions are functorial, we need to choose such decompositions functorially as well. This can be achieved for any profinite set by inspecting the proof of implication $(2)\Rightarrow (1)$ in \cite[\href{https://stacks.math.columbia.edu/tag/08ZY}{Tag 08ZY}]{stacks-project} and observing that the posited diagram is functorial in $S$. We will however simply write $S=\limit_i S_i$ and assume the mentioned construction tacitly.
\end{remark}

If $A$ is finitely generated, the definition of $A_\blacksquare$ is analogous to $\Z_\blacksquare$ from example \ref{example: examples for pre-analytic rings}, otherwise one has to approximate $A_\blacksquare$ from below, using finitely generated subrings. 
\begin{definition}[The pre-analytic ring $A_\blacksquare$ for finitely generated $A$]
	\label{definition: the pre-analytic ring Ablack for finitely generated A}
	
	Let $A$ be a discrete finitely generated $\Z$-algebra. The pre-analytic ring $A_\blacksquare$ has underlying ring $A$, free completion $S=\limit_i S_i \mapsto \limit_i \free{A}{S_i}$ and corresponding natural transformation induced by the natural maps $S_i\to \free{A}{S_i}$ for finite $S_i$.
\end{definition}

To define $A_\blacksquare$ in full generality we will need a relative version of the previous definition.
\begin{definition}[The pre-analytic ring $(A, R)_\blacksquare$ for finitely generated $R$]
	Let $\phi\colon R\to A$ be a map of discrete $\Z$-algebras, such that $R$ is finitely generated. The pre-analytic ring $(A, R)_\blacksquare$ has underlying ring $A$, free completion $\scalarExtension{\phi}\free{R_\blacksquare}{-} = A\otimes_R \free{R_\blacksquare}{-}$ and corresponding natural transformation induced by the morphisms $S\to \free{R_\blacksquare}{S}$.
\end{definition}

We can now define $A_\blacksquare$ in full generality.
\begin{definition}[The pre-analytic ring $A_\blacksquare$]
	\label{definition: the pre-analytic ring Ablack}
	
	Let $A$ be a discrete ring. For any extremally disconnected $S$ define
		$$\free{A_\blacksquare}{S}\ldef \colimit_{A'\subseteq A} \free{(A, A')_\blacksquare}{S},$$
	where the colimit is filtered and runs over all finitely generated $\Z$-subalgebras $A'$ of $A$. Now $A_\blacksquare$ is the pre-analytic ring with underlying ring $A$, free completion given by $\free{A_\blacksquare}{-}$ and corresponding natural transformation induced by any of the compatible morphisms $S\to \free{(A, A')_\blacksquare}{S}$.
\end{definition}

\begin{remark}[$A_\blacksquare$ as a colimit of pre-analytic rings.]
	In definition \ref{definition: the pre-analytic ring Ablack}, $A_\blacksquare$ can be viewed as the colimit of the pre-analytic rings $(A, A')_\blacksquare$ along the various $A'$, all with the same underlying ring $A$. We will however not really need this characterization and thus opted for the more explicit definition.
\end{remark}

\begin{remark}[Sanity check]
	Suppose that in definition \ref{definition: the pre-analytic ring Ablack} the $\Z$-algebra $A$ is a finitely generated. Then clearly $A'\ldef A$ is cofinal in the partially ordered indexing set of the colimit, ensuring that the natural map
		$$\free{A_\blacksquare^{\text{old}}} S \cong \colimit_{A' = A} \free{A'_\blacksquare} S \to \colimit_{A'\subseteq A} \free{A'_\blacksquare} S \rdef \free{A_\blacksquare} S$$
	is an isomorphism, recovering the definition \ref{definition: the pre-analytic ring Ablack for finitely generated A} of the finitely generated case.
\end{remark}

We can now reintroduce the relative ring $(A, R)_\blacksquare$ for a general base $R$.
\begin{definition}[The pre-analytic ring $(A, R)_\blacksquare$]
	Let $\phi\colon R\to A$ be a morphism of discrete rings. Then $(A, R)_\blacksquare$ is the pre-analytic ring with underlying ring $A$, with free completion $\scalarExtension{\phi}\free{R_\blacksquare}{-} = A\otimes_R \free{R_\blacksquare}{-}$ and corresponding natural transformation induced by the maps $S\to \free{R_\blacksquare}{S}$.
\end{definition}

\begin{remark}[Free completions for profinite sets]
	\label{remark: free completions for profinite sets}
	
	Observe, that for all the pre-analytic rings introduced in this section, the free completions can readily be generalized to an arbitrary profinite set $S$. Furthermore, as a consequence of \cite[Prop. 5.6]{condenseddotpdf} we observe that this generalized definition agrees with the construction mentioned in remark \ref{remark: when is the tensor product the derived tensor product?}. In particular, the complex constructed there is always concentrated in degree $0$ and hence $\Dotimes_\preAnaRing{A}$ really is the left derived functor of $\otimes_\preAnaRing{A}$ for call cases of $\preAnaRing{A}$ relevant to us -- of course only once we know the rings $A_\blacksquare$ and $(A, R)_\blacksquare$ to be analytic.
\end{remark}

\begin{remark}[Solidness]
	These (pre-)analytic rings $A_\blacksquare$ play an essential role in much of the algebraic theory of condensed mathematics, so much so, that $A_\blacksquare$-complete modules deserve a special name. A condensed $A$-module is called \emph{$A$-solid} if it is $A_\blacksquare$-complete. A $\Z$-solid condensed abelian group is simply called \emph{solid}. We will use this terminology occasionally as well.
\end{remark}

Let us now try to work towards an understanding of what \emph{being solid} means for a module. Recall that \emph{solidness} is characterized by existence of certain linear maps from the condensed modules $\free{A_\blacksquare}{S}$ with $S$ extremally disconnected. It is thus a good idea to develop an understanding of these $A_\blacksquare$-free modules themselves. We will try to do so for even for a general profinite set $S$.
\begin{definition}[The abelian group of continuous maps]
	For any profinite set $S$ and any discrete abelian group $M$ denote by $\continuous{S}{M}$ the discrete topological abelian group given by pointwise addition.
\end{definition}

We can now prove a first result regarding the structure of $A_\blacksquare$-free abelian groups.
\begin{lemma}[Finitely-free $A_\blacksquare$-complete modules as $A$-duals]
	\label{lemma: finitely-free A-solid modules as A-duals}
	
	Let $A$ be a discrete finitely generated $\Z$-algebra. For any profinite $S$ there is a natural isomorphism
		$$\free{A_\blacksquare}{S} \cong \associated{\intHom_\Ab(\continuous S \Z, A)}\cong \intHom_\condensedAb(\associated{\continuous{S}{\Z}}, \associated{A})$$
	of $A$-modules.
\end{lemma}
\begin{proof}
	Write $S\cong \limit_i S_i$ where each $S_i$ is finite and discrete. Then $\continuous{S_i}{\Z}\cong \bigoplus_{S_i} \Z$ is free of finite rank and hence there is a natural isomorphism
		$$\free {A} {S_i}\cong \intHom_\Ab(\continuous{S_i}{\Z}, A).$$
	As $\free{A_\blacksquare}{S} \ldef \limit_i \free{A}{S_i}$, we in particular obtain
		$$\free{A_\blacksquare}{S}\cong \limit_i \associated{\intHom_\Ab(\continuous{S_i}{\Z}, A)}\cong \associated{\intHom_\Ab(\continuous{S}{\Z}, A)}.$$
	This shows the first isomorphism. As both $\continuous S \Z$ and $A$ are discrete, additive maps $\continuous S \Z\to A$ as abelian groups and as topological abelian groups agree, so the second isomorphism follows as well.
\end{proof}

Thus, these $A_\blacksquare$-free modules are $A$-duals of abelian groups of continuous maps, at least for finitely generated $A$. One can interpret elements of this dual space as a kind of measure.
\begin{definition}[Pseudo-measures]
	Let $S$ be a set and $\mathcal{G}$ be a $\cup$-stable, $\cap$-stable and complement-stable system of subsets of $S$ such that $\emptyset\in \mathcal{G}$. Let $A$ be a discrete abelian group. An \emph{$A$-valued pseudo-measure} on $S$ is a function
		$$\mu\colon \mathcal{G} \to A$$
	such that $\mu(\emptyset)=0$ and $\mu(U\cup V) = \mu(U) + \mu(V)$ for any two disjoint $U, V\in\mathcal{G}$. Tacitly assuming the system of sets to be fixed, we denote by $\measures{S}{A}$ the set of pseudo-measures on $S$.
\end{definition}

If $S$ is profinite, one such system of subsets is given by the clopen subsets of $S$. Indeed, $\emptyset$ is clopen, the union and intersection of clopen subsets is clopen and certainly the complement of a clopen set is clopen as well. Furthermore, by identifying a clopen subset with its indicator function we obtain the following characterization.
\begin{lemma}
	\label{lemma: indicator functions generate}
	
	Let $S$ be a profinite set. Then
		$$\continuous{S}{\Z} = \langle 1_U\colon S\to \Z\setseparator U\subseteq S\text{ clopen}\rangle_\Z$$
	that is any continuous function $S\to \Z$ is a finite $\Z$-linear combination of indicator functions of clopen subsets of $S$.
\end{lemma}
\begin{proof}
	If $U\subseteq S$ is clopen then $1_U\colon S\to \Z$ is continuous. Indeed, if $V$ is a(n open) subset of $\Z$ then its preimage is either $\emptyset$, the clopen $U$, its complement $S\setminus U$ or the whole space $S$ -- all of which are open in $S$. Certainly, also any finite $\Z$-linear combination of such indicator functions is continuous as well.
	
	For the converse assume that $f\colon S\to \Z$ is continuous. As $S$ is compact so is its image $f(S)$ and since $\Z$ is discrete so is its subspace $f(S)$. But a discrete space is compact if and only if it is finite hence $f$ takes only finitely many values. Let $n\in f(S)$. Then $\{n\}$ is clopen in $\Z$ (by discreteness of $\Z$) and hence so is the fiber $f^{-1}(\{n\})$. It follows that
		$$f = \sum_{n\in f(S)} n \cdot 1_{f^{-1}(\{n\})}$$
	is such a $\Z$-linear combination as claimed.
\end{proof}

We can thus characterize duals of such groups $\continuous{S}{\Z}$.
\begin{lemma}[Finitely-free $A_\blacksquare$-modules as spaces of measures]
	Let $A$ be a discrete finitely generated $\Z$-algebra. For any profinite set $S$ there is a natural isomorphism of $A$-modules
		$$\intHom_\Ab(\continuous S \Z, A) \cong \measures S A$$
	between the $A$-dual of $\continuous S \Z$ and the $A$-valued pseudo-measures on the system of clopen subsets of $S$. In particular, lemma \ref{lemma: finitely-free A-solid modules as A-duals} yields that
		$$\free{A_\blacksquare}{S} \cong \associated{\measures S A}$$
	is the condensed set associated to the space of $A$-valued pseudo-measures on $S$.
\end{lemma}
\begin{proof}
	Recall that by the previous lemma \ref{lemma: indicator functions generate} the abelian group $\continuous{S}{\Z}$ is generated by indicator functions $1_U$ of clopen subsets $U$ of $S$. As stated above, the result then follows from identifying these clopen subsets $U$ of $S$ with their indicator functions $1_U$. More precisely, map an $f\colon \continuous S \Z \to A$ to the pseudo-measure $\mu_f\colon U\mapsto f(1_U)$ and map every pseudo-measure $\mu$ to the additive function $f_\mu$ such that $1_U\mapsto \mu(U)$ for every clopen subset $U$ of $S$. Clearly $\mu_{-}$ and $f_{-}$ are additive. Furthermore, they satisfy
		$$\mu_{f_\mu} = \mu\text{ and }f_{\mu_f}=f$$
	so are inverse of each other.
\end{proof}

With this interpretation of $\Z_\blacksquare$-free abelian groups at hand, we can try to understand solidness.
\begin{remark}[Interpreting $\Z_\blacksquare$-completeness]
	\label{remark: interpreting Zblack-completeness}
	
	Suppose that the $\Z$-module $\associated{M}$ is $\Z_\blacksquare$-complete. If $f\colon S\to M$ is a map from a profinite set, then we obtain by completeness a commuting triangle
	$$\begin{tikzcd}[column sep=1pt]
		S\ar[rr, "f"]\ar[dr] &&M\\
		&\measures{S}{\Z}\ar[ur, "\hat f", dotted, swap]
	\end{tikzcd}$$
	Now if $\mu\in \measures S \Z$ is any pseudo-measure, then we can define $\int f\d \mu\ldef \hat f(\mu)$. Since the map $S\to\measures{S}{\Z}$ corresponds to mapping any $s \in S$ to its corresponding Dirac-pseudo-measure $\delta_s \in \measures{S}{\Z}$, this gives the familiar identity $\int f \d\delta_s = \hat f(\delta_s) = f(s)$. Furthermore, since $f\mapsto \hat f$ is $\Z$ linear we obtain further that $\int (\alpha f + \beta g) \d\mu = \alpha \int f\d\mu + \beta \int g\d\mu$ for all $f, g\colon S\to M$ and integers $\alpha, \beta$. We find that this \emph{integral} really has the expected properties. Thus, $\Z_\blacksquare$-completeness of $M$ can be characterized as saying that \emph{any continuous function $S\to M$ is integrable against any $\Z$-valued pseudo-measure on $M$}.
	
%	In fact, since the $\Z$-submodule of $\measures{S}{\Z}$ generated by all Dirac-measures is dense in $\measures{S}{\Z}$ (with the weak-$\ast$-topology), it necessarily follows that for any $f\colon S\to M$ and any $\mu\in\measures{S}{\Z}$ the integral $\int f\d\mu$ is in fact the Lebesgue integral (for any $f$ both maps $\mu \mapsto \int f \d\mu$ are continuous and agree on the dense subset spanned by Dirac-measures). Thus $\Z_\blacksquare$-completeness of $M$ can be characterized as saying that \emph{any continuous function $S\to M$ is integrable against any $\Z$-valued measure on $M$}.
	
	Let us investigate this integrability in the special case of the profinite set $S = \hat\N=\N\cup\{\infty\}$, that is the one-point-compactification of $\N$. Then a continuous map $f\colon \hat \N\to M$ corresponds precisely to a convergent sequence $(a_n)_{n\in \N}$ in $M$ with limit $a_\infty = f(\infty)$. Denoting by $\continuousZero{\N}{M}$ the null-sequences in $M$, we obtain a decomposition $\continuous {\hat \N} M \cong \continuousZero {\N} M \oplus M$ of abelian groups by mapping $f\mapsto (f\vert_\N-f(\infty), f(\infty))$. This induces a decomposition of pseudo-measures $\measures{\hat\N}{\Z} \cong \measures{\N}{\Z}\oplus \Z$, by means of $\mu\mapsto (\mu\vert_\N, \mu(\infty))$. As $\N$ is a discrete space, any subset is clopen hence the system of clopen subsets of $\N$ agrees with the usual $\sigma$-algebra (both are simply the powerset of $\N$). It is clear that any measure on $\N$ is a pseudo-measure on $\N$ (but not conversely as $\sigma$-additivity is not clear). Now, using the linearity of the integral we see that for any $f\in \continuous{S}{\Z}$ and any $\mu \in \measures{S}{\Z}$ we have
		$$\int_{\hat \N} f \d\mu = \int_{\hat \N} f-f(\infty) \d\mu + f(\infty)\mu(\infty).$$
%	But $(f-f(\infty))(\infty) = 0$, yielding further that $\int_{\hat\N} f-f(\infty)\d\mu = \int_\N f - f(\infty)\d\mu$ and hence that
%		$$\int_{\hat \N} f \d\mu = \int_\N f-f(\infty)\d\mu + f(\infty)\mu(\infty) = \quote{\sum_{n\in \N} (f(n) - f(\infty))\mu(n)} + f(\infty)\mu(\infty).$$
	In particular if $f(\infty) = 0$ then there is an element
		$$\sum_{n\in \N} f(n) \mu(n) \ldef \int_{\hat\N} f\d\mu \in M,$$
	that is \emph{any null-sequence in $M$ is summable against any $\Z$-valued pseudo-measure $\mu$ on $\hat \N$}. By choosing a pseudo-measure $\mu$ with $\mu(n)=1$ for all $n\in \N$ we see that for any null-sequence $f$ in $M$ the sum $\sum_{n\in \N} f(n)$ ought to exist in $M$ -- it follows that solidness is a \emph{non-archimedean} notion. Indeed, in an archimedean setting, this is surely false. \Eg in $\R$ one has the standard example of the (partial sums of the) harmonic series which diverge.
\end{remark}

A rather important result on the structure of $\continuous S \Z$, due to Nöbling in \cite{noebling}, based on results of Specker in \cite{specker}, is given in the following theorem. We will only sketch the proof (using an argument attributed to Bergman in \cite[Thm. 97.2]{fuchs-infinite-abelian-groups}) -- a formally verified proof can be found in \cite{asgeirsson}.
\begin{theorem}[Specker's theorem]
	\label{theorem: specker's theorem}
	
	For any profinite set $S$, the discrete abelian group $\continuous{S}{\Z}$ of continuous integer valued functions on $S$ is free with rank at most $2^{\vert S\vert}$.
	
	In particular, there is some set $I$ with $\vert I\vert\le 2^{\vert S\vert}$ such that for any discrete finitely generated $\Z$-algebra we have an isomorphism of $A$-modules
		$$\free{A_\blacksquare}{S} \cong \prod_I \associated{A}.$$
\end{theorem}
\begin{proofsketch}
	Pick a continuous injection $S\to \prod_\lambda \{0, 1\}$ onto a closed subset, where $\lambda$ is suitably large and assumed an ordinal (here one can take $\lambda$ to be equicardinal to the set of clopen subsets of $S$ such that for each clopen subset the corresponding projection to $\{0, 1\}$ is the indicator function on the subset). In particular, the elements $\mu \in \lambda$ correspond to ordinals $\mu < \lambda$. For such a $\mu < \lambda$ consider the idempotent $e_\mu \in \continuous S \Z$ given by $S\to\prod_\lambda \{0, 1\}\to \{0, 1\}\subseteq \Z$ induced by the projection onto the $\mu$-th factor. For $r\in \N_0$, order the products $e_{\mu_1}\cdot (\cdots) \cdot e_{\mu_r}$ with $\mu_1 > \dots > \mu_r$ lexicographically. Denote by $E$ the set of such products that cannot be written as a linear combination of smaller such products. We claim that $E$ is a basis of $\continuous{S}{\Z}$. We will proceed by induction on the ordinal $\lambda$. Denote by $\tilde S$ the image of $S$ and for any $\mu <\lambda$ denote by $\tilde S_\mu$ the image of $S$ in $\prod_{\mu'<\mu} \{0, 1\}$ after projecting from $\prod_\lambda \{0, 1\}$. 
	
	If $\lambda = 0$ then $S=\{\ast\}$ and the claim is trivial.
	
	If $\lambda$ is a limit ordinal, \ie $\lambda = \bigcup_{\mu < \lambda} \mu$, then $\tilde S \cong \colimit_{\mu < \lambda} \tilde S_\mu$ is the union of the $\tilde S_\mu$. Denoting by $E_\mu$ the analogous subsets of $\continuous{\tilde S_\mu}{\Z}$ -- which are bases by induction, we observe that $E \cong \colimit_{\mu < \lambda} E_\mu$ is a basis for $\continuous{\tilde S}{\Z}$.
	
	Finally assume that $\lambda = \rho + 1$ is a successor ordinal. Then the above map identifies with the map $S \hookrightarrow \tilde S_\rho \times \{0, 1\}$. Denote by $\tilde S_0$ and $\tilde S_1$ the intersections $\tilde S\cap (\tilde S_\rho \times \{0\})$ and $\tilde S\cap (\tilde S_\rho \times \{1\})$ respectively. As intersections of closed subsets, these are closed. Furthermore, we have $\tilde S = \tilde S_0 \cup \tilde S_1$. Denote by $\tilde S'$ the intersection of the images of $\tilde S_0$ and $\tilde S_1$ under the projection $\tilde S_\rho\times\{0, 1\}\to \tilde S_\rho$. We thus have the short exact sequence
		$$0\to \continuous{\tilde S_\rho}\Z\to \continuous {\tilde S} \Z \to \continuous{\tilde S'}{\Z}\to 0$$
	where the left morphism maps $f\colon \tilde S_\rho\to\Z$ to the continuous function $(s, b)\mapsto f(s)$, and the right morphism maps $g\colon \tilde S\to \Z$ to the difference of restrictions $s\mapsto f(s,0) - f(s, 1)$. Now applying the induction hypothesis for both $\tilde S_\rho$ and $\tilde S'$ (mapping to $\prod_{\mu<\rho} \{0, 1\}$) we see that the basis vectors of $E$ that do not start with $e_\rho$ form a basis of $\continuous{\tilde S_\rho}{\Z}$, while the basis vectors that do start with $e_\rho$ project to a basis of $\continuous{\tilde S'}{\Z}$. So $E$ really is a basis by exactness of the above sequence.
	
	For the second claim follows by choosing an isomorphism $\continuous{S}{\Z}\cong \bigoplus_I \Z$ of abelian groups. Then
		$$\free{A_\blacksquare}S \cong \associated{\intHom_\Ab(\continuous{S}{\Z}, A)} \cong \associated{\intHom_\Ab(\bigoplus\nolimits_I \Z, A)}=\prod_I \associated{A},$$
	as $\intHom_\Ab(-, A)$ maps direct sums to products and as $X\mapsto\associated{X}$ commutes with them.
\end{proofsketch}

To get an intuition on how $A_\blacksquare$-completed tensor products behave, we state the following proposition but simply reference its proof.
\begin{proposition}
	\label{proposition: tensor products of compacts, f.g. case}
	
	Let $A$ be a discrete finitely generated $\Z$-algebra. If $M=\prod_I \associated{A}$ and $N=\prod_J \associated{A}$, then
	$$M\Dotimes_{A_\blacksquare} N \cong \prod_{I\times J} \associated{A}.$$
\end{proposition}
\begin{shortproof}
	In case of $A=\Z$ this is \cite[Prop. 6.3]{condenseddotpdf}. The other case is analogous.
\end{shortproof}

A sample consequence of this proposition is given in
\begin{example}[Some completed tensor products]
	One can calculate the following tensor products directly:
	\begin{itemize}
		\item $\Z_p \Dotimes_{\Z_\blacksquare} \R \cong 0$
		\item $\Z_p \Dotimes_{\Z_\blacksquare} \Z_\ell \cong 0$
		\item $\Z_p \Dotimes_{\Z_\blacksquare} \Z_p \cong \Z_p$
		\item $\Z_p \Dotimes_{\Z_\blacksquare} \series{\Z}{T} \cong \series{\Z_p}{T}$
		\item $\series{\Z}{U} \Dotimes_{\Z_\blacksquare} \series{\Z}{T} \cong \series{\Z}{U, T}$
	\end{itemize}
	for primes $p\ne \ell$. The precise reasoning is given in \cite[Example 6.4]{condenseddotpdf}. What one should observe is that the (derived) completed tensor product $\Dotimes_{\Z_\blacksquare}$ \quote{seems to ask both factors for what $I$-adic topology they have, and then combines} them. In particular solidness and more generally $A$-solidness for some ring $A$ are non-archimedean in nature (see also remark \ref{remark: interpreting Zblack-completeness}). Clausen and Scholze however state, that solidness is very well suited for applications in $p$-adic functional analysis. If one wishes to obtain a good condensed foundation for real or complex functional analysis, one has to put in significantly more work. This is the content of the (now completed, but very subtle) \emph{Liquid Tensor Experiment}, see \href{https://xenaproject.wordpress.com/2020/12/05/liquid-tensor-experiment/}{for a short introduction}. 
\end{example}

This allows us to calculate tensor products even in the general case. Recall that in case of an $A$ that is not finitely generated we have $\free{A_\blacksquare}{S} = \colimit_{A'\subseteq A} A\otimes_{A'}\free{A'_\blacksquare}{S}$ for any extremally disconnected $S$. With $I$ as in Specker's theorem \ref{theorem: specker's theorem}, we thus obtain that $\free{A_\blacksquare}{S} \cong \colimit_{A'\subseteq A} A\otimes_{A'} \prod_I A'$. First, observe that all objects of this form are already compact projective.
\begin{lemma}[More compact projectives]
	\label{lemma: more compact projectives}
	
	Let $A$ be a discrete ring such that $A_\blacksquare$ is analytic. Then for any set $I$ the $A$-module
		$$\colimit_{A'\subseteq A} A\otimes_{A'} \prod_I A'$$
	is $A_\blacksquare$-complete and a compact projective object of $\mod{A_\blacksquare}$.
\end{lemma}
\begin{proof}
	Choose a large enough extremally disconnected $S$ such that the set $J$ associated to it by Specker's theorem \ref{theorem: specker's theorem} satisfies $\vert I\vert \le \vert J \vert$. Choosing an injection $I \hookrightarrow J$ we find that for all $A'$ the product $\prod_{I} A'$ is a retract of the product $\prod_J A'$ where the retraction is natural. Hence, the $A$-module in question is a retract of
		$$\colimit_{A'\subseteq A} A\otimes_{A'} \prod_J A'\cong \free{A_\blacksquare}{S}$$
	which is compact projective. As being $A_\blacksquare$-complete and being compact projective are stable under retracts it follows that $\colimit_{A'\subseteq A} A\otimes_{A'} \prod_I A'$ is $A_\blacksquare$-complete and compact projective as well.
\end{proof}

\begin{corollary}
	\label{corollary: tensor products of compacts, general case}
		
	Let $A$ be a discrete ring such that $A_\blacksquare$ is analytic. If $M = \colimit_{A'\subseteq A} A\otimes_{A'} \prod_I A'$ and $M = \colimit_{A'\subseteq A} A\otimes_{A'} \prod_j A'$ for two sets $I$ and $J$, then
		$$M\Dotimes_{A_\blacksquare} N \cong M\otimes_{A_\blacksquare} N\cong \colimit_{A'\subseteq A} A\otimes_{A'} \prod_{I\times J} A'.$$
	In particular, by the previous lemma \ref{lemma: more compact projectives} the tensor product $M\otimes_{A_\blacksquare} N$ is compact projective in $\mod{A_\blacksquare}$. Furthermore, $M\Dotimes_{A_\blacksquare} N$ is derived compact in $\derived{A_\blacksquare}$.
\end{corollary}
\begin{proof}
	By this previous lemma, we also know $M$ and $N$ to be projective. We immediately obtain that
		$$M\Dotimes_{A_\blacksquare} N \cong M\otimes_{A_\blacksquare} N,$$
	and we do not need to work in the derived setting. Since $A_\blacksquare$ is analytic, we know by theorem \ref{theorem: properties of the category of complete modules} that $A_\blacksquare$-completion exists and is strong monoidal, yielding that
		$$M\otimes_{A_\blacksquare} N \cong \completion{\mleft(\colimit_{A'} A\otimes_{A'} \prod_I A'\mright)\otimes_A \mleft(\colimit_{A''} A\otimes_{A''} \prod_J A''\mright)}{A_\blacksquare}.$$
	Since tensor products commute with arbitrary colimits we find that
	\begin{align*}
		M\otimes_{A_\blacksquare} N &\cong \completion{\colimit_{A'}\colimit_{A''} \mleft(A\otimes_{A'} \prod_I A'\mright) \otimes_A \mleft(A\otimes_{A''} \prod_J A''\mright)}{A_\blacksquare}\\
		&\cong \completion{\colimit_{A'}\colimit_{A''} \mleft(\prod_I A'\mright)\otimes_{A'} A \otimes_A A\otimes_{A''} \mleft(\prod_J A''\mright)}{A_\blacksquare}\\
		&\cong \completion{\colimit_{A'}\colimit_{A''} \mleft(\prod_I A'\mright)\otimes_{A'} A \otimes_{A''} \mleft(\prod_J A''\mright)}{A_\blacksquare}.
	\end{align*}
	By injecting the pair $(A', A'')$ into $(A''', A''')$ where $A''' \ldef \Z[A'\cup A'']$ is $\Z$-finitely generated, we find that the diagonal $A'=A''$ is cofinal and hence that
	\begin{align*}
		M\otimes_{A_\blacksquare} N &\cong \completion{\colimit_{A'} \mleft(\prod_I A'\mright)\otimes_{A'} A \otimes_{A'} \mleft(\prod_J A'\mright)}{A_\blacksquare}\\
		&\cong \completion{\colimit_{A'} A\otimes_{A'}\mleft(\prod_I A'\mright)\otimes_{A'}\mleft(\prod_J A'\mright)}{A_\blacksquare}.
	\end{align*}
	Now $\completion{-}{A_\blacksquare}$ is a left adjoint so commutes with arbitrary colimits. Furthermore, if we denote by $\phi\colon A'\to A$ the inclusion then $\completion{-}{A_\blacksquare} \circ \scalarExtension{\phi} \cong \completedScalarExtension{\phi}{A_\blacksquare} \circ \completion{-}{A'_\blacksquare}$ as follows from the proof of proposition \ref{proposition: scalar restriction and extension of complete modules}, \ie $A_\blacksquare \otimes_{A} A\otimes_{A'} - \cong A_\blacksquare\otimes_{A'} - \cong A_\blacksquare\otimes_{A'_\blacksquare} A'_\blacksquare\otimes_{A'} -$. Thus,
	\begin{align*}
		M\otimes_{A_\blacksquare} N &\cong \colimit_{A'} A_\blacksquare\otimes_{A'_\blacksquare}\completion{\mleft(\prod_I A'\mright)\otimes_{A'}\mleft(\prod_J A'\mright)}{A'_\blacksquare}\\
		 &\cong \colimit_{A'} A_\blacksquare\otimes_{A'_\blacksquare}\mleft(\prod_I A'\mright)\otimes_{A'_\blacksquare}\mleft(\prod_J A'\mright).
	\end{align*}
	By proposition \ref{proposition: tensor products of compacts, f.g. case} we thus have
		$$M\otimes_{A_\blacksquare} N \cong \colimit_{A'} A_\blacksquare\otimes_{A'_\blacksquare} \prod_{I\times J} A' \cong \completion{\colimit_{A'} A\otimes_{A'}\prod_{I\times J} A'}{A_\blacksquare}$$
	which is simply
		$$\colimit_{A'} A\otimes_{A'} \prod_{I\times J} A'$$
	as this $A$-module is already $A_\blacksquare$-complete by the previous lemma.
\end{proof}

These previous statements assert that the tensor product of compact objects is still compact. We obtain the following result as a consequence that occasionally proves useful.
\begin{lemma}[Solid derived compacts are enriched compact]
	\label{lemma: solid derived compacts are enriched compact}
	
	Let $A$ be a discrete ring such that $A_\blacksquare$ is analytic. Let $K$ be a derived compact object of $\derived{A_\blacksquare}$. Then $K$ is $\condensedAb$-enriched compact.
	
	Similarly, if $R\to A$ is a morphism of discrete rings such that $(A, R)_\blacksquare$ is analytic, then any derived compact of $\derived{(A, R)_\blacksquare}$ is $\condensedAb$-enriched compact.
\end{lemma}
\begin{proofsketch}
	By theorem \ref{theorem: properties of the category of complete modules} we know that if $Y$ is $A_\blacksquare$-complete and $S$ is extremally disconnected then $\eRHom_A(\free{A_\blacksquare}{S}, Y) \cong \eRHom_A(\free{A}{S}, Y)$ in $\derived{\condensedAb}$. This implies that also $\intRHom_A(\free{A_\blacksquare}{S}, Y) \cong \intRHom_A(\free{A}{S}, Y)$ as $\eRHom_A$ is the underlying condensed abelian group of $\intRHom_A$. By applying $H^0$ we thus obtain that $\intHom_A(\free{A_\blacksquare}{S}, Y) \cong \intHom_A(\free{A}{S})$ in $\mod{A}$. Using this, we can show that $\mod{A_\blacksquare}$ is actually closed monoidal. Indeed, we find for all $X, Y\in\mod{A_\blacksquare}$ the $A$-module $\intHom_A(X, Y)$ to be $A_\blacksquare$-complete. For this let $S$ be extremally disconnected. Then
	\begin{align*}
		\Hom(\free{A_\blacksquare}{S}, \intHom_A(X, Y)) &\cong \Hom(X, \intHom_R(\free{A_\blacksquare}{S}, Y))\\
		& \cong \Hom(X, \intHom_A(\free{A}{S}, Y)) \cong \Hom(\free{A}{S}, \intHom_A(X, Y))
	\end{align*}
	showing the desired isomorphy. One can then check that this really is a partial right adjoint to $\otimes_{A_\blacksquare}$.
	
	Now knowing that $\mod{A_\blacksquare}$ is closed monoidal we find by corollary \ref{corollary: a sufficient condition for enriched compactness} that if $-\Dotimes_{A_\blacksquare} K$ preserves compactness then $K$ is $\mod{A_\blacksquare}$-enriched compact. But the tensor products of the derived compact generators of $\derived{A_\blacksquare}$ are again derived compact by corollary \ref{corollary: tensor products of compacts, general case} so this follows. This shows $\mod{A_\blacksquare}$-enriched compactness of $K$. Now as the $\condensedAb$-enrichment is simply given by the underlying condensed abelian groups of the $\mod{A_\blacksquare}$-enrichment it follows that $K$ is also already $\condensedAb$-enriched compact (since direct sums in all categories in question agree).
	
	The case of $(A, R)_\blacksquare$ is analogous, after one observes that the relevant claim about tensor products of generators being compact again translates.
\end{proofsketch}

Having now investigated the rings $(A, R)_\blacksquare$ and $A_\blacksquare$ on their own, we should then set them in relation to each other. The following lemma provides various natural maps between them.
\begin{lemma}[Some maps of pre-analytic rings]
	\label{lemma: some maps of pre-analytic rings}
	
	We have the following maps of pre-analytic rings
	\begin{enumerate}[(i)]
		\item
		For any discrete ring $A$ a map $A\to A_\blacksquare$ with underlying morphism of rings $\id_A$.
		
		\item
		\label{lemma: some maps of pre-analytic rings, ii}
		For any morphism $\phi\colon R\to A$ of discrete rings a map $R_\blacksquare\to A_\blacksquare$ with morphism of underlying rings $\phi$.
		
		\item
		For any morphism $R\to A$ of discrete rings a map $(A, R)_\blacksquare\to A_\blacksquare$ with morphism of underlying rings $\id_A$.
		
		\item
		For any commutative square
		$$\begin{tikzcd}
			R\ar[r]\ar[d] &A\ar[d, "\phi"]\\
			R'\ar[r] & A'
		\end{tikzcd}$$
		of discrete rings a map $(A, R)_\blacksquare \to (A', R')_\blacksquare$ with morphism of underlying rings $\phi$.
	\end{enumerate}
\end{lemma}
\begin{proof}
	We will only show the existence of the desired natural morphisms between (finitely) free modules and omit the compatibility with the natural transformations of underlying condensed sets.
	\begin{enumerate}[(i)]
		\item
		Suppose that the $\Z$-algebra $A$ is finitely generated and that $S \cong \limit_i S_i$. Then the projections $S\to S_i$ induce natural morphisms $\free{A}{S}\to \free{A}{S_i}$ which in turn induce a natural morphism $\free{A}{S} \to \limit_i \free{A}{S_i} \rdef \free{A_\blacksquare}{S}$, giving the desired $A\to A_\blacksquare$.
		
		Now suppose that $A$ is general. If $A'$ is any finitely generated $\Z$-subalgebra of $A$ then by the previous argument we have a morphism $A'\to A'_\blacksquare$. These induce $A$-linear maps
		$$\free{A}{S} \cong \colimit_{A'\subseteq A} \free{A'}{S}\cong \colimit_{A'\subseteq A} A\otimes_{A'} \free{A'}{S} \to \colimit_{A'\subseteq A} A\otimes_{A'}\free{A'_\blacksquare}{S}$$
		natural in the extremally disconnected $S$, giving the desired $A\to A_\blacksquare$.
		
		\item
		Suppose that the $\Z$-algebras $R$ and $A$ are finitely generated. If $S\cong \limit_i S_i$ then the natural maps $\free{R}{S_i} \to \free{A}{S_i}$ induce a natural $R$-linear map $\free{R_\blacksquare}{S} \ldef \limit_i \free{R}{S_i} \to \limit_i \free{A}{S_i} \rdef \free{A_\blacksquare}{S}$, giving the desired $R_\blacksquare\to A_\blacksquare$.
		
		Now suppose that $A$ is general, but $R$ is still finitely generated. Then the poset of $\Z$-finitely generated $R$-subalgebras $A''$ of $A$ is cofinal in the poset of all finitely generated $\Z$-subalgebras of $A$, as any $\Z$-subalgebra $A'$ can be embedded in the still $\Z$-finitely generated $R$-algebra $A''=\Z[A', R]$. Thus, the natural morphism of colimits
		$$\colimit_{R\to A''\subseteq A} A\otimes_{A''}\free{A''_\blacksquare}{S} \to \colimit_{A'\subseteq A} A\otimes_{A'}\free{A'_\blacksquare}{S}\rdef \free{A_\blacksquare}{S}$$
		is an isomorphism. Then by the previous paragraph we have compatible $R$-linear maps $\free{R_\blacksquare}{S} \to \free{A''_\blacksquare}{S} \to A\otimes_{A''} \free{A''_\blacksquare}{S}$ that then induce a natural $R$-linear map into the first, hence into the second colimit which is $\free{A_\blacksquare}{S}$, giving the desired $R_\blacksquare \to A_\blacksquare$.
		
		Now if both $R$ and $A$ are general, then any finitely generated $\Z$-subalgebra $R'$ gives rise to a map $R'_\blacksquare \to A_\blacksquare$ by the previous paragraph. Hence we obtain an $R$-linear map
		$$\free{R_\blacksquare}{S} \ldef \colimit_{R'\subseteq R} R\otimes_{R'}\free{R'_\blacksquare}{S} \to \free{A_\blacksquare}{S}$$
		natural in the extremally disconnected $S$ from the compatible morphisms $R\otimes_{R'} \free{R'_\blacksquare}{S} \to R\otimes_{R'} \free{A_\blacksquare}{S} \to \free{A_\blacksquare}{S}$ where the second map is $R$-multiplication on the $R$-module $\free{A_\blacksquare}{S}$. This gives the desired $R_\blacksquare \to A_\blacksquare$.
		
		\item
		By \ref{lemma: some maps of pre-analytic rings, ii} we have a map $R_\blacksquare\to A_\blacksquare$. For any extremally disconnected $S$ we thus obtain the $A$-linear map
		$$\free{(A, R)_\blacksquare}{S}\ldef A\otimes_R \free{R_\blacksquare}{S} \to A\otimes_R \free{A_\blacksquare}{S} \to \free{A_\blacksquare}{S}$$
		natural in $S$.
		
		\item 
		By \ref{lemma: some maps of pre-analytic rings, ii} there is a map $R_\blacksquare \to R'_\blacksquare$. Denote its corresponding natural transformation $\free{R_\blacksquare}{-}\Rightarrow \free{R'_\blacksquare}{-}$ by $\eta$. Observe that the morphism $R\to R'$ induces a natural transformation $\otimes_R \Rightarrow \otimes_{R'}$. Thus we obtain $A$-linear maps
		$$\free{(A, R)_\blacksquare}{S} \ldef A\otimes_R \free{R_\blacksquare}{S} \xrightarrow{\phi\otimes \eta_S} A'\otimes_R \free{R'_\blacksquare}{S} \to A'\otimes_{R'} \free{R'_\blacksquare}{S} \rdef \free{(A', R')}{S}$$
		natural in the extremally disconnected $S$, giving the desired $(A, R)_\blacksquare \to (A', R')_\blacksquare$.
	\end{enumerate}
\end{proof}



\newpage
\section{Analyticity of \texorpdfstring{$A_\blacksquare$ and $(A, R)_\blacksquare$}{Solid Rings}}
\label{section: analyticity of Ablack and (A, R)black}

The goal of this chapter is to prove that the pre-analytic rings $A_\blacksquare$ and $(A, R)_\blacksquare$ are analytic. The case of $(A, R)_\blacksquare$ follows essentially formally from knowledge of analyticity of $R_\blacksquare$ by lemma \ref{lemma: partial analyticity of (A, R)black}. Thus the main difficulty lies in $A_\blacksquare$ and we will obtain the result by tackling the analyticity of $A_\blacksquare$ for increasingly more general $A$, the most basic instance being $\Z_\blacksquare$.
\begin{theorem}[Analyticity of $\Z_\blacksquare$]
	\label{theorem: analyticity of Zblack}
	
	The pre-analytic ring $\Z_\blacksquare$ is analytic.
\end{theorem}
%\begin{proof}
%	Due to Specker's theorem \ref{theorem: specker's theorem} any finitely $\Z_\blacksquare$-free module $\free{\Z_\blacksquare}T$ for some extremally disconnected $T$ is of the form $\prod_I \associated{\Z}$ for some set $I$. Instead of considering a complex $C^\bullet = (\dots \to C^{-1} \to C^0 \to 0)$ where each $C^i$ is a $\Z_\blacksquare$-free modules, allow more generally that each $C^i$ is a direct sums of products of the form $\prod_I \associated{\Z}$ for various sets $I$.
%	
%	We will show that for each extremally disconnected set $S$ the induced morphism
%	$$\eRHom_\Z(\free{\Z_\blacksquare}{S}, C^\bullet)\to \eRHom_\Z(\free{\Z}{S}, C^\bullet)$$
%	is an isomorphism, implying analyticity.
%	
%	Again by Specker's theorem \ref{theorem: specker's theorem} there is some set $I$ such that $\continuous S \Z \cong \bigoplus_I \Z$. An immediate consequence was that $\free{\Z_\blacksquare}{S}$ is the condensed abelian group associated to the space of measures $\intHom_\Ab(\continuous{S}{\Z}, \Z)\cong \measures{S}{\Z}$. Consider analogously $\measures{S}{\R} \cong \intHom_\Ab(\continuous{S}{\Z}, \R)$ and $\measures{S}{\R/\Z} \cong \intHom_\Ab(\continuous{S}{\Z}, \R/\Z)$, the spaces of $\R$-valued and $\R/\Z$-valued measures on $S$ respectively. Akin to $\measures S \Z \cong \prod_I \Z$ we then have $\measures S \R \cong \prod_I \R$ and $\measures{S}{\R/\Z}\cong \prod_I \R/\Z$. Now $\continuous{S}{\Z}\cong\bigoplus_I \Z$ is free, hence projective, so the image
%	$$0\to \measures S \Z \to \measures{S}{\R} \to \measures S {\R/\Z}\to 0$$
%	of $0\to \Z\to\R\to \R/\Z\to 0$ under $\intHom_\Ab(\continuous{S}{\Z}, -)$ is exact.
%	
%	\begin{todo}
%		finish this
%	\end{todo}
%\end{proof}
While the claim of theorem \ref{theorem: analyticity of Zblack} is very interesting and fundamental to the rest of this thesis (and hence deserves a full proof), the argument given in \cite{condenseddotpdf} by Clausen and Scholze is a rather long and technical calculation. Hence, it will simply be referenced.
\begin{shortproof}[of theorem \ref{theorem: analyticity of Zblack}]
	This is shown in \cite[Lecture VI, pp. 38--41]{condenseddotpdf}
\end{shortproof}


The next level of generality will be analogous to Hilbert's basis theorem.
\begin{theorem}[A kind of Hilbert's basis theorem]
	\label{theorem: a kind of hilbert's basis theorem}
	
	If $R$ is a discrete finitely generated $\Z$-algebra such that $R_\blacksquare$ is analytic, then the pre-analytic ring $R[T]_\blacksquare$ is analytic as well.
\end{theorem}
The proof of this theorem uses many intermediate results which will be needed again at a later time, hence we will devote to this proof its own section \ref{section: an analogue of hilbert's basis theorem}. Assuming the statement proven for now, we immediately obtain the following corollary.
\begin{corollary}[Analyticity of $A_\blacksquare$ for finitely generated free $\Z$-algebras]
	\label{corollary: analyticity of Ablack for finitely generated free Z-algebras}
	
	The pre-analytic ring $\Z[T_1,\dots, T_n]_\blacksquare$ associated to the polynomial ring in $n$ variables over $\Z$ is analytic.
\end{corollary}
\begin{shortproof}
	This follows inductively from theorem \ref{theorem: a kind of hilbert's basis theorem} and theorem \ref{theorem: analyticity of Zblack}.
\end{shortproof}

The next stepping stone will be analyticity of any finitely generated $\Z$-algebra, that is precisely the quotients of the previous rings $\Z[T_1,\dots, T_n]$. The idea is essentially to choose an epimorphism $R\ldef\Z[T_1, \dots, T_n] \twoheadrightarrow A$ and observe that $A_\blacksquare\cong (A, R)_\blacksquare$. Then one obtains analyticity of $A_\blacksquare$ after noticing that analyticity of $(A, R)_\blacksquare$ can always be reduced to $R_\blacksquare$ as shown in lemma \ref{lemma: partial analyticity of (A, R)black}. To prove this reduction, one needs to control the derived scalar extension of (finitely) $R_\blacksquare$-free modules. Clausen and Scholze show this in the case that $A$ and $R$ are both finitely generated $\Z$-algebras in \cite[Intro. Appendix to Lec. VIII]{condenseddotpdf}. In this thesis, we will need the generalization where we do not require $R$ or $A$ to be finitely generated.
\begin{lemma}[Derived scalar extension is concentrated in degree $0$]
	\label{lemma: derived scalar extension is concentrated in degree 0}
	
	Let $\phi\colon R\to A$ be a morphism of discrete rings. Then for any extremally disconnected set $S$ the natural morphism
		$$\LeftD\scalarExtension{\phi}\free{R_\blacksquare}{S}\ldef A\Dotimes_R \free{R_\blacksquare} S\to A\otimes_R \free{R_\blacksquare}S \rdef \scalarExtension{\phi}\free{R_\blacksquare}{S}$$
	in $\derived{\associated{A}}$ is an isomorphism.
\end{lemma}
\begin{proof}
	We show the analogous claim for the (derived) tensor product with an arbitrary $R$-module $M$ instead of $A$. Observe that an $A$-linear map is an isomorphism if and only if it is an isomorphism as an $R$-linear map. We will thus show that the given morphism is an isomorphism of $R$-modules.
	
	Assume that $R$ is a finitely generated $\Z$-algebra. Then by Specker's theorem \ref{theorem: specker's theorem} there is a set $I$ for which $\free{R_\blacksquare}{S}\cong \prod_I R$. Now $R$ is Noetherian and thus coherent, so by \cite[\href{https://stacks.math.columbia.edu/tag/05CZ}{Tag 05CZ}]{stacks-project} the product $\prod_I R$ and hence $\free{R_\blacksquare}{S}$ are $R$-flat. In particular the natural morphism
		$$M\Dotimes_R \free{R_\blacksquare}{S} \to M\otimes_R \free{R_\blacksquare}{S}$$
	is a (quasi-)isomorphism.
	
	Now assume that $R$ is general. Recall that tensor products (as left adjoints) commute with colimits. As colimits of complexes are calculated degree-wise we observe further that also $\otimes_\category{A}^\bullet$ commutes with colimits of complexes. We thus obtain (by replacing $M\Dotimes_\category{A} -$ by $P^\bullet\otimes_\category{A}^\bullet -$ where $P^\bullet$ is a suitable left resolution of $M$) that
		$$M\Dotimes_R \free{R_\blacksquare}{S} \cong \colimit_{R'\subseteq R} M \Dotimes_R (R \otimes_{R'} \free{R'_\blacksquare}{S})$$
	where the colimit is calculated in chain complexes (but considered in $\derived{\associated{R}}$), is filtered and runs over all finitely generated $\Z$-subalgebras $R'$ of $R$. But by the previous case we already know that $R\Dotimes_{R'}\free{R'_\blacksquare}{S} \to R\otimes_{R'}\free{R'_\blacksquare}{S}$ is a (quasi-)isomorphism and hence we obtain by exactness of $\colimit_{R'\subseteq R}$ that
		$$\colimit_{R'\subseteq R} M \Dotimes_R (R \otimes_{R'} \free{R'_\blacksquare}{S})\cong \colimit_{R'\subseteq R} M \Dotimes_R R \Dotimes_{R'} \free{R'_\blacksquare}{S}\cong \colimit_{R'\subseteq R} M \Dotimes_{R'} \free{R'_\blacksquare}{S}.$$
	But again by the previous case (and exactness of $\colimit_{R'\subseteq R}$) we have
		$$\colimit_{R'\subseteq R} M\Dotimes_{R'} \free{R'_\blacksquare}{S} \isorightarrow \colimit_{R'\subseteq R} M\otimes_{R'} \free{R'_\blacksquare}{S}\cong \colimit_{R'\subseteq R} M\otimes_R R\otimes_{R'} \free{R'_\blacksquare}{S} \cong M\otimes_R \free{R_\blacksquare}{S},$$
	concluding the proof.
	
%	Choose $I$ as in Specker's theorem \ref{theorem: specker's theorem}, so that $\free{R_\blacksquare}{S} \cong \prod_I R$. Both $- \otimes_R \free{R_\blacksquare} S$ and $- \Dotimes_R \free{R_\blacksquare} S$ commute with filtered colimits, hence assume without loss of generality that $M$ is a finitely generated $R$-module. Since $R$ is Noetherian, $M$ thus admits a resolution $\epsilon\colon F^\bullet\to M$ from a complex of finite free $R$-modules. Hence
%		$$M\Dotimes_R \free{R_\blacksquare}{S}\cong F^\bullet\otimes_R \prod_I R \isorightarrow \prod_I F^\bullet \cong \prod_I M$$
%	in $\derived{\associated{R}}$, where the second isomorphism holds holds on the level of complexes since each $F^i$ is finite free, so the induced morphism $F^i\otimes_R \prod_I R \to \prod_I F^i \otimes_R R \cong \prod_I F^i$ of $R$-modules is an isomorphism, as the tensor product commutes with finite direct sums. But by \cite[\href{https://stacks.math.columbia.edu/tag/059K}{Tag 059K}]{stacks-project} we have
%		$$M\otimes_R \free{R_\blacksquare}{S} \cong M\otimes_R \prod_I R \isorightarrow \prod_I M\otimes_R R \cong \prod_I M$$
%	as well, since $M$ is finitely presented.
\end{proof}

We need one further technical result for lemma \ref{lemma: partial analyticity of (A, R)black}.
\begin{lemma}[Relative completeness]
	\label{lemma: relative completeness}
	
	
	Let $R\to A$ be a morphism of discrete rings such that $R_\blacksquare$ is analytic. Then the scalar restriction of any $(A, R)_\blacksquare$-free module is $R_\blacksquare$-complete.
\end{lemma}
\begin{proof}
	Since scalar restriction commutes with direct sums, we can assume with out of loss of generality that $M$ is of the form $\free{(A, R)_\blacksquare}{S}$ for some extremally disconnected set $S$. Observe that $\scalarRestriction{\phi}\free{(A, R)_\blacksquare}{S} = A\otimes_R \free{R_\blacksquare}{S}$. We will show more generally that for any discrete $R$-module $N$ the tensor product $N\otimes_R \free{R_\blacksquare}{S}$ is complete. For this choose an $R$-presentation $\bigoplus_I R \to \bigoplus_J R\to N \to 0$ of $N$. As $-\otimes_R \free{R_\blacksquare}{S}$ commutes with arbitrary colimits and is right exact, the induced sequence
		$$\bigoplus\nolimits_I \free{R_\blacksquare}{S} \to \bigoplus\nolimits_J \free{R_\blacksquare}{S}\to N \otimes_R \free{R_\blacksquare}{S} \to 0$$
	is exact. Thus $N\otimes_R \free{R_\blacksquare}{S}$ is $R_\blacksquare$-presentable hence $R_\blacksquare$-complete.
\end{proof}

We can now show the claimed partial result on analyticity of $(A, R)_\blacksquare$.
\begin{lemma}[Partial analyticity of $(A, R)_\blacksquare$]
	\label{lemma: partial analyticity of (A, R)black}
	
	Let $\phi\colon R\to A$ be a morphism of discrete rings such that $R_\blacksquare$ is analytic. Then $(A, R)_\blacksquare$ is analytic as well.
\end{lemma}
\begin{proof}
	Let $S$ be any extremally disconnected set and $C^\bullet$ a complex of $(A, R)_\blacksquare$-free modules, concentrated in non-positive degrees. By lemma \ref{lemma: derived scalar extension is concentrated in degree 0} we obtain $\scalarExtension{\phi}\free{R_\blacksquare} S \cong \LeftD\scalarExtension{\phi}\free{R_\blacksquare}S$. As $\scalarExtension{\phi}\ladj \scalarRestriction{\phi}$ and as $\scalarRestriction{\phi}$ is exact we obtain by lemma \ref{lemma: enriched adjunctions are stable under derivation} that in $\derived{\condensedAb}$ we have
	\begin{align*}
		\eRHom_A(\free{(A,R)_\blacksquare}{S}, C^\bullet)
		&= \eRHom_A(\scalarExtension{\phi}\free{R_\blacksquare}{S}, C^\bullet)\\
		&\cong \eRHom_A(\LeftD\scalarExtension{\phi}\free{R_\blacksquare}{S}, C^\bullet)\\
		&\cong\eRHom_R(\free{R_\blacksquare}{S}, \RightD\scalarRestriction{\phi}C^\bullet) \cong \eRHom_R(\free{R_\blacksquare}{S},\scalarRestriction{\phi}C^\bullet)
	\end{align*}
	where the last isomorphism follows since $\RightD\scalarRestriction{\phi} = \scalarRestriction{\phi}$ by exactness of $\scalarRestriction{\phi}$. But the complex $\scalarRestriction{\phi}C^\bullet$ is $R_\blacksquare$-complete by lemma \ref{lemma: relative completeness} and we thus further obtain that
		$$\eRHom_R(\free{R_\blacksquare}{S},\scalarRestriction{\phi}C^\bullet) \cong \eRHom_R(\free R S, \scalarRestriction{\phi}C^\bullet)\cong \eRHom_A(\free A S, C^\bullet)$$
	since clearly $\LeftD\scalarExtension{\phi}\free R S \cong \scalarExtension{\phi}\free R S \cong \free A S$ as $\free R S$ is projective.
\end{proof}

%The final ingredient for proving analyticity for any finitely generated $\Z$-algebra is now
%\begin{lemma}[The ring $(A, R)_\blacksquare$ for well-behaved extensions]
%	\label{lemma: the ring (A, R)solid for well-behaved extensions}
%	
%	\begin{todo}
%		Need to generalize this to not hold only for finitely generated rings
%	\end{todo}
%	
%	If $\phi\colon R\to A$ is a finite-type integral extension of discrete rings such that $R$ is Noetherian. Then for any set $I$ the natural morphism
%		$$\scalarExtension{\phi}\prod_I R = A \otimes_R \prod_I R \to \prod_I A$$
%	of $A$-modules is an isomorphism.
%	
%	In particular if both $R$ and $A$ are finitely generated $\Z$-algebras, then for any extremally disconnected $S$ the morphism $\free{(A, R)_\blacksquare}{S} \to \free{A_\blacksquare}{S}$ is an isomorphism, as it is of the above form by Specker's theorem \ref{theorem: specker's theorem}. Thus the map of (pre-)analytic rings $(A, R)_\blacksquare\to A_\blacksquare$ is an isomorphism.
%\end{lemma}
%\begin{proof}
%	By assumption the $R$-module $A$ is generated by finitely many integral elements. As $R$ is Noetherian, $A$ is even finitely presented. Hence by \cite[\href{https://stacks.math.columbia.edu/tag/059K}{Tag 059K}]{stacks-project} the natural morphism
%	$$A\otimes_R\prod_I R \to \prod_I A \otimes_R R \cong \prod_I A$$
%	of $R$-modules is an isomorphism. This proves the claim, since this map is $A$-linear as well.
%\end{proof}

The final ingredient for proving analyticity for any finitely generated $\Z$-algebra is now to show that $(A, R)_\blacksquare$ and $A_\blacksquare$ are isomorphic as pre-analytic rings in case that $R\to A$ is a finite morphism. Clausen and Scholze argue this in \eg \cite[Exposition to Lecture IX]{condenseddotpdf} in case of finitely generated $\Z$-algebras. Again, we will need the generalization to arbitrary rings later on. We already state and prove the generalization here.
\begin{lemma}[Analytic rings associated to finite morphisms]
	\label{lemma: analytic rings associated to finite morphisms}
	
	If $\phi\colon R\to A$ is a finite morphism of discrete rings, then the morphism of pre-analytic rings $(A, R)_\blacksquare \to A_\blacksquare$ with underlying map of rings $\id_A$ from \ref{lemma: some maps of pre-analytic rings} is an isomorphism.
\end{lemma}
\begin{proof}
	Consider first the case that both $R$ and $A$ are finitely-generated $\Z$-algebras. Let $S$ be extremally disconnected and choose $I$ as in Specker's theorem \ref{theorem: specker's theorem}. Since $R\to A$ is finite, the $R$-module $A$ is a finitely generated (and hence finitely presented) and \cite[\href{https://stacks.math.columbia.edu/tag/059K}{Tag 059K}]{stacks-project} asserts that the natural morphism
		$$A\otimes_R \prod_I R \to \prod_I A\otimes_R R \cong \prod_I A$$
	is an isomorphism. But again by Specker, this is precisely the morphism $\free{(A, R)_\blacksquare}{S} \to \free{A_\blacksquare}{S}$, proving that $(A, R)_\blacksquare \to A_\blacksquare$ is an isomorphism (with inverse given by the inverse natural transformation $\free{A_\blacksquare}{-}\Rightarrow \free{(A, R)_\blacksquare}{-}$ and with the same underlying morphism $\id_A$).
	
	Now assume that $R$ and $A$ are arbitrary. Consider again some extremally disconnected $S$. Consider the partially ordered sets $P_R$ and $P_A$ of $\Z$-finitely generated subalgebras of $R$ and $A$ respectively. Denote by $P$ the category of finite morphisms $R'\to A'$ with $R' \in P_R$ and $A' \in P_A$ such that
	$$\begin{tikzcd}
		R' \ar[r]\ar[d] &A'\ar[d]\\
		R \ar[r] & A
	\end{tikzcd}$$
	commutes. By \cite[\href{https://stacks.math.columbia.edu/tag/0BTG}{Tag 0BTG}]{stacks-project} the two projections $P\to P_R$ and $P\to P_A$ are final. Then
	\begin{align*}
		\free{(A, R)_\blacksquare}{S}
		&\ldef A\otimes_R (\colimit_{R'} R\otimes_{R'} \free{R'_\blacksquare}{S})\\
		&\cong \colimit_{R'} A\otimes_{R'} \free{R'_\blacksquare}{S}\\
		&\cong \colimit_{R'\to A'}A\otimes_{R'} \free{R'_\blacksquare}{S}\\
		&\cong \colimit_{R'\to A'}A\otimes_{A'} A' \otimes_{R'} \free{R'_\blacksquare}{S}\\
		&\cong \colimit_{R'\to A'} A\otimes_{A'} \free{A'_\blacksquare}{S}\\
		&\cong \colimit_{A'} A\otimes_{A'} \free{A'_\blacksquare}{S} \rdef \free{A_\blacksquare}{S}
	\end{align*}
	where the first isomorphism is due to the tensor product $\otimes_R$ commuting with filtered colimits, the second and fifth isomorphism being due to finalness of $P\to P_R$ and $P\to P_A$ respectively, and the fourth isomorphism following from the previous case applied to each $R'\to A'$. By naturality of all isomorphisms it follows that $\free{(A, R)_\blacksquare}{S}\to\free{A_\blacksquare}{S}$ is an isomorphism as desired.
\end{proof}

\begin{theorem}[Analyticity of finite type $\Z$-algebras]
	\label{theorem: analyticity of finite type Z-algebras}
	
	For every finitely generated discrete $\Z$-algebra $A$, the ring $A_\blacksquare$ is analytic.
\end{theorem}
\begin{proof}
	Choose a surjection $R\ldef \Z[T_1, \dots, T_n]\twoheadrightarrow A$. Then $R_\blacksquare$ is analytic by corollary \ref{corollary: analyticity of Ablack for finitely generated free Z-algebras} and hence lemma \ref{lemma: partial analyticity of (A, R)black} gives the analyticity of $(A, R)_\blacksquare$. But lemma \ref{lemma: analytic rings associated to finite morphisms} yields that $A_\blacksquare \cong (A, R)_\blacksquare$ as any surjection of rings is finite (indeed a morphism of rings is surjective if and only if it is finite and epic by \cite[\href{https://stacks.math.columbia.edu/tag/04VT}{Tag 04VT}]{stacks-project}). Hence, $A_\blacksquare$ is analytic.
%	\begin{todo}
%		check out https://arxiv.org/pdf/2105.12591, Prop.2.11, 3.15, 3.16, 3.17
%	\end{todo}
\end{proof}

We immediately obtain the following corollary.
\begin{corollary}[Analyticity of $(A, R)_\blacksquare$ for finitely generated $R$]
	\label{corollary: analyticity of (A, R)black for finitely generated R}
	
	For every map $R\to A$ of discrete $\Z$-algebras for which $R$ is finitely generated, the ring $(A, R)_\blacksquare$ is analytic.
\end{corollary}
\begin{shortproof}
	This immediately follows from proposition \ref{theorem: analyticity of finite type Z-algebras} together with lemma \ref{lemma: partial analyticity of (A, R)black}.
\end{shortproof}


This allows us to prove the final stage, namely that $A_\blacksquare$ is analytic for any discrete ring $A$. While Clausen and Scholze hint at the fact that this should be true (\eg immediately after \cite[Definition 9.1]{condenseddotpdf} -- albeit in slightly different form), they do not provide an argument in their notes.
\begin{theorem}[Analyticity of $A_\blacksquare$ for non-finite $\Z$-algebras $A$]
	\label{theorem: analyticity of Ablack for non-finite Z-algebras A}
	
	For any ring discrete $A$ the pre-analytic ring $A_\blacksquare$ is analytic.
\end{theorem}
\begin{proof}	
	Let $A'\subseteq A$ be a finitely generated subring. Then the pre-analytic ring $(A, A')_\blacksquare$ is analytic by corollary \ref{corollary: analyticity of (A, R)black for finitely generated R}.
	%	where $\eta_S$ is the scalar extension of $\free{A'}{S}\to \free{A'_\blacksquare}{S}$ after identifying $\scalarExtension{\phi}\free{A'}{S}=\free{A}{S}$. Suppose that $C^\bullet$ is a complex of $\preAnaRing{A}_{A'}$-free $A$-modules. Observe that
	%		$$\free{\preAnaRing{A}_{A'}}{S} \ldef \scalarRestriction{\phi}\scalarExtension{\phi}\free{A'_\blacksquare}{S} = A\otimes_{A'} \free{A'_\blacksquare}{S}$$
	%	is $A'_\blacksquare$-free since $A$ is discrete and $A'_\blacksquare$ is analytic by assumption. Thus $\scalarRestriction{\phi}C^\bullet$ is a complex of $A'_\blacksquare$-complete modules and we obtain that
	%	\begin{align*}
		%		\eRHom_A(\scalarExtension{\phi}\free{A'_\blacksquare}{S}, C^\bullet) &\cong \eRHom_{A'}(\free{A'_\blacksquare}{S}, \scalarRestriction{\phi}C^\bullet)\\
		%		&\cong \eRHom_{A'}(\free{A'}{S}, \scalarRestriction{\phi}C^\bullet)\cong \eRHom_A(\free{A}{S}, C^\bullet)
		%	\end{align*}
	%	since $A'_\blacksquare$ is analytic by assumption, implying that the pre-analytic ring $\scalarExtension{\phi}A'_\blacksquare$ is analytic.
	
	
	Consider the poset $P\ldef \{A''\subseteq A\setseparator A'' \text{ finitely generated}\}$ of such finitely generated subrings of $A$ and its sub-poset $P'\ldef \{A'' \in P\setseparator A'\subseteq A''\}$ of those subrings that contain $A'$. Denote be $i\colon P'\to P$ the inclusion. Observe that $i$ is a final functor. Indeed, for any given $A''\in P$ the comma category $(A''/i)=\{ A''' \in P'\setseparator A''\subseteq A'''\} = A''/P'$ is connected: This category is non-empty since the subring $\Z[A' \cup A''] \in P'$ contains $A''$. Furthermore, any two $A'''$ and $A''''$ are connected by a zig-zag $A'''\subseteq \Z[A'''\cup A''''] \supseteq A''''$. 
	
	In particular for any extremally disconnected set $S$ we obtain that the natural morphism of $A$-modules
	$$\colimit_{A'\subseteq A''\subseteq A} \free{(A, A'')_\blacksquare}{S} \to \colimit_{A''\subseteq A} \free{(A, A'')_\blacksquare}{S} \rdef \free{A_\blacksquare}{S}$$
	is an isomorphism. 
	
	Now for any such $A'\subseteq A''$ both $(A, A')_\blacksquare$ and $(A, A'')_\blacksquare$ are analytic by \ref{corollary: analyticity of (A, R)black for finitely generated R}. Thus, by proposition \ref{proposition: maps of analytic rings from maps of pre-analytic rings} the existence of the map of pre-analytic rings $(A, A')_\blacksquare \to (A, A'')_\blacksquare$ from lemma \ref{lemma: some maps of pre-analytic rings} implies that $\id_A$ is a map of analytic rings $(A, A')_\blacksquare \to (A, A'')_\blacksquare$. Thus, by proposition \ref{proposition: scalar restriction and extension of complete modules} each $\free{(A, A'')_\blacksquare}{S}$ is already $(A, A')_\blacksquare$-complete, implying that $\free{A_\blacksquare}{S} = \colimit_{A'\subseteq A'' \subseteq A} \free{(A, A'')_\blacksquare}{S}$ as a colimit of $(A, A')_\blacksquare$-complete modules is $(A, A')_\blacksquare$-complete as well. By a similar argumentation we find that the inclusion $A'\subseteq A$ is a map of analytic rings $A'_\blacksquare \to (A, A')_\blacksquare$, further implying that the scalar restriction of $\free{A_\blacksquare}{S}$ to $A'$ is also $A'_\blacksquare$-complete, by proposition \ref{proposition: scalar restriction and extension of complete modules}.
	
	Now suppose that $C^\bullet$ is a complex of $A_\blacksquare$-free modules concentrated in non-positive degrees. Then we have just shown that $C^\bullet$ is $(A, A')_\blacksquare$-complete and its scalar restriction to $A'$ is $A'_\blacksquare$-complete for any $A'\in P$. Now for some extremally disconnected $S$ we have
	\begin{align*}
		&\eRHom_A(\free{A_\blacksquare}{S}, C^\bullet) \\
		=\,&\eRHom_A(\colimit\nolimits_{\phi\colon A'\subseteq A} \free{(A, A')_\blacksquare}{S}, C^\bullet)\\
		\cong\,&\homotopyLimit\nolimits_{\phi\colon A'\subseteq A} \eRHom_A(\free{(A, A')_\blacksquare}{S}, C^\bullet)
	\end{align*}
	by \assumptionPreWord \ref{assumption: right derived Hom maps colimits to limits}. By lemma \ref{lemma: derived scalar extension is concentrated in degree 0} the natural map $\LeftD\scalarExtension{\phi}\free{A'_\blacksquare}{S}\to \scalarExtension{\phi}\free{A'_\blacksquare}{S}$ is a quasi-isomorphism. Recall that $\scalarExtension{\phi}\ladj \scalarRestriction{\phi}$ and that $\scalarRestriction{\phi}$ is exact. Hence, we obtain from lemma \ref{lemma: enriched adjunctions are stable under derivation} that in $\derived{\condensedAb}$ we have
	\begin{align*}
		\eRHom_A(\free{(A, A')_\blacksquare}{S}, C^\bullet)&= \eRHom_A(\scalarExtension{\phi}\free{A'_\blacksquare}{S}, C^\bullet)\\
		& \cong \eRHom_A(\LeftD\scalarExtension{\phi}\free{A'_\blacksquare}{S}, C^\bullet)\\
		&\cong \eRHom_{A'}(\free{A'_\blacksquare}{S}, \RightD\scalarRestriction{\phi}C^\bullet)\\
		&\cong \eRHom_{A'}(\free{A'_\blacksquare}{S}, \scalarRestriction{\phi}C^\bullet)\\
		&\cong \eRHom_{A'}(\free{A'}{S}, \scalarRestriction{\phi}C^\bullet) \cong \eRHom_A(\free{A}{S}, C^\bullet)
	\end{align*}
	since $\RightD\scalarRestriction{\phi} \cong \scalarRestriction{\phi}$, since $\scalarRestriction{\phi}C^\bullet$ is $A'_\blacksquare$-complete and since $\LeftD\scalarExtension{\phi}\free{A'}{S} = \free{A}{S}$. Applying this to the previous chain of isomorphisms we obtain that
	$$\eRHom_A(\free{A_\blacksquare}{S}, C^\bullet) \cong \homotopyLimit\nolimits_{\phi\colon A'\subseteq A} \eRHom_A(\free{A}{S}, C^\bullet) \cong \eRHom_A(\free{A}{S}, C^\bullet)$$
	by another application of \assumptionPreWord \ref{assumption: right derived Hom maps colimits to limits}. Thus, $A_\blacksquare$ is analytic as well.
\end{proof}

We finally obtain the relative case in full generality.
\begin{corollary}[Analyticity of $(A, R)_\blacksquare$]
	\label{corollary: analyticity of (A, R)black}
	
	For any morphism of discrete rings $R\to A$ the pre-analytic ring $(A, R)_\blacksquare$ is analytic.
\end{corollary}
\begin{shortproof}
	By theorem \ref{theorem: analyticity of Ablack for non-finite Z-algebras A} the ring $R_\blacksquare$ is analytic, hence by lemma \ref{lemma: partial analyticity of (A, R)black} so is $(A, R)_\blacksquare$.
\end{shortproof}


We can now translate the morphism of pre-analytic rings from lemma \ref{lemma: some maps of pre-analytic rings} into the analytic world.
\begin{corollary}[Some maps of analytic rings]
	\label{corollary: some maps of analytic rings}
	
	We have the following maps of analytic rings
	\begin{enumerate}[(i)]
		\item
		For any discrete ring $A$ the map $\id_A\colon A\to A_\blacksquare$.
		
		\item
		Any morphism $\phi\colon R\to A$ of discrete rings is the map $\phi\colon R_\blacksquare\to A_\blacksquare$.
		
		\item
		For any morphism $R\to A$ of discrete rings the map $\id_A\colon (A, R)_\blacksquare\to A_\blacksquare$.
		
		\item
		For any commutative square
			$$\begin{tikzcd}
				R\ar[r]\ar[d] &A\ar[d, "\phi"]\\
				R'\ar[r] & A'
			\end{tikzcd}$$
		of discrete rings the map $\phi\colon (A, R)_\blacksquare \to (A', R')_\blacksquare$.
	\end{enumerate}
\end{corollary}
\begin{proof}
	By proposition \ref{proposition: maps of analytic rings from maps of pre-analytic rings} it is enough to exhibit in each of the cases a map of pre-analytic rings with matching morphism of underlying rings, as this map of underlying rings then is a map of analytic rings. But lemma \ref{lemma: some maps of pre-analytic rings} provides exactly such maps.
\end{proof}


\newpage
\section{An Analogue of Hilbert's Basis Theorem}
\label{section: an analogue of hilbert's basis theorem}

The goal of this section is provide a proof of theorem \ref{theorem: a kind of hilbert's basis theorem} -- the analogue of Hilbert's basis theorem.

\begin{convention}
	Throughout this section let $R$ be a finitely generated discrete $\Z$-algebra. Denote by $A\ldef R[T]$ the ring of polynomials and by $A_\infty\ldef \laurent{R}{T^{-1}}$ the ring of \emph{formal Laurent series in $T^{-1}$} (\ie series $\sum_{n\gg \infty} r_n T^{-n}$ with finite principal part), with its natural $\ideal{T^{-1}}$-adic topology.
	
	Assume that $R_\blacksquare$ is analytic. Then by lemma \ref{lemma: partial analyticity of (A, R)black} also $(A, R)_\blacksquare$ is analytic and the identity $\id_A$ is a morphism of analytic rings $R_\blacksquare \to (A, R)_\blacksquare$ by corollary \ref{corollary: some maps of analytic rings}, implying that any $(A, R)_\blacksquare$-complete module is also $R_\blacksquare$-complete by proposition \ref{proposition: scalar restriction and extension of complete modules}.
\end{convention}

\begin{lemma}
	\label{lemma: Ainfty is A-flat, f.g. case}
	
	The $A$-module $A_\infty$ is $A$-flat.
\end{lemma}
\begin{proof}
	Observe that $A_\infty = \laurent{R}{T^{-1}} = \series{R}{T^{-1}}[T]$ is the $\ideal{T^{-1}}$-adic completion of the ring $A'\ldef R[T, T^{-1}]$. Since $R$ is $\Z$-finitely generated and hence Noetherian, Hilbert's basis theorem asserts that $A'\cong R[T, U]/\ideal{UT - 1}$ is Noetherian as well. Thus by \cite[\href{https://stacks.math.columbia.edu/tag/00HT}{Tag 00HT}]{stacks-project}, the completion is $A'$-flat. But $A' = R[T, T^{-1}]\cong R[T]_T = A_T$ is a localization of $A$, hence $A$-flat by \cite[\href{https://stacks.math.columbia.edu/tag/00HT}{Tag 00HT}]{stacks-project}. As flatness is transitive by \cite[\href{https://stacks.math.columbia.edu/tag/00HT}{Tag 00HT}]{stacks-project}, we obtain that $A_\infty$ is $A$-flat.
\end{proof}

\begin{lemma}
	\label{lemma: resolution for Ainfty, f.g. case}
	There is an exact sequence of $A$-modules
		$$0\to A \otimes_R \series R U \xrightarrow{UT - 1} A \otimes_R \series R U \to A_\infty\to 0$$
	where the first two terms are compact projectives of $\mod{(A, R)_\blacksquare}$. In particular $A_\infty$ is $(A, R)_\blacksquare$-complete, $R_\blacksquare$-complete and quasi-isomorphic to a bounded complex of compact projectives of $\mod{(A, R)_\blacksquare}$ hence compact in $\derived{(A, R)_\blacksquare}$.
\end{lemma}
\begin{proof}
	Observe that $A \otimes_R \series R U \cong \mleft(\bigoplus_\N R\mright)\otimes_R \series{R}{U}\cong \bigoplus_\N \series{R}{U}\cong \series R U [T]$ is the ring of polynomials over the ring of power series, hence exactness of the sequence immediately follows as $A_\infty = \laurent{R}{T^{-1}} = \series{R}{T^{-1}}[T]$. It remains to show that $A\otimes_R \series{R}{U}$ is a compact projective object of $\mod{(A, R)_\blacksquare}$. But by (the proof of) lemma \ref{lemma: more compact projectives} the object $A\otimes_R \series{R}{U} \cong A\otimes_R \prod_\N R$ is a compact projective of $\mod{(A, R)_\blacksquare}$ (as the retraction is preserved under $A\otimes_R -$).
	
%	Denote by $\hat \N$ the one-point-compactification of $\N$ and observe that the sequence of $R$-modules
%		$$0\to \free{R_\blacksquare}{\infty}\to\free{R_\blacksquare}{\hat \N} \to \series{R}{U} \to 0$$
%	where the first morphism is induced by the inclusion $\{\infty\}\to \hat\N$ and where the second morphism corresponds to the map $\hat \N\to\series{R}{U}$ given by $n\mapsto U^n$ and $\infty\mapsto 0$, is exact. Denote the inclusion $\free{R_\blacksquare}{\infty} \to \free{R_\blacksquare}{\hat\N}$ by $s$ and denote by $r\colon \free{R_\blacksquare}{\hat\N}\to \free{R_\blacksquare}{\infty}$ the morphism induced by the projection $\hat\N\to\{\infty\}$. By functoriality of $\free{R_\blacksquare}{-}$ we obtain that $r\circ s = \id$ as $\{\infty\}\to \hat\N\to\{\infty\}$ is the identity. In particular $r$ is a retraction for $s$ and the short exact sequence is split. As any additive functor preserves split exact sequences, the induced sequence
%		$$0\to A\otimes_R \free{R_\blacksquare}{\infty}\to A\otimes_R \free{R_\blacksquare}{\hat\N}\to A\otimes_R \series{R}{U}\to 0$$
%	is split exact as well, implying that $A\otimes_R\series{R}{U}$ is a retract of $\free{(A, R)_\blacksquare}{\hat\N} = A\otimes_R \free{R_\blacksquare}{\hat \N}$. But the latter is projective in $\mod{(A, R)_\blacksquare}$ and so must be its retract. As for any complex $C^\bullet\in\derived{(A, R)_\blacksquare}$ the functor $\Hom(-, C^\bullet)\colon \derived{(A, R)_\blacksquare}\to \Ab$ is cohomological by \cite[\href{https://stacks.math.columbia.edu/tag/0149}{Tag 0149}]{stacks-project} and the first two terms of the short exact sequence are compact, an application of the $5$-lemma yields, that $A\otimes_R\series{R}{U}$ is compact as well.
\end{proof}


We will need the following statement even for general rings $R$ that are not $\Z$-finitely generated.
\begin{lemma}[Formal Laurent series are idempotent]
	\label{lemma: formal Laurent series are idempotent}
	
	If $R$ is any ring then multiplication 
	$$\laurent{R}{T^{-1}}\otimes_{R[T]}\laurent{R}{T^{-1}}\to \laurent{R}{T^{-1}}$$
is an $R[T]$-module isomorphism.
\end{lemma}
\begin{proof}
	Define a morphism in the other direction such that $a\mapsto 1\otimes a$. It is clear that the composition $\laurent{R}{T^{-1}} \to \laurent{R}{T^{-1}}\otimes_{R[T]} \laurent{R}{T^{-1}}\to \laurent{R}{T^{-1}}$ is the identity. To show that the other composition is the identity as well, we need that $1\otimes ab = a\otimes b$ for all $a, b\in \laurent{R}{T^{-1}}$. For this write $a = \sum_{k\gg -\infty} r_k T^{-k}$ giving a sequence $(a_n)_{n\in \N}$ with $a_n \ldef \sum_{n \ge k \gg -\infty} r_k T^{-k}$ polynomials converging to $a$ in the $\ideal{T^{-1}}$-adic topology. Write further $a_n = T^{v_n} \hat{a_n}$ with $v_n\in \Z$ and $\hat{a_n} \in R[T]$ not divisible by $T$. Observe that since the tensor product is $T$-balanced and hence also $T^{-1}$-balanced. Indeed for $x,y \in \laurent{R}{T^{-1}}$ arbitrary we find
	$$xT^{-1}\otimes y = xT^{-1}\otimes TT^{-1}y = xT^{-1}T\otimes T^{-1}y = x\otimes T^{-1}y$$
	as $T^{-1} \in \laurent{R}{T^{-1}}$ exists. We now find that
	$$1\otimes a_nb = 1\otimes T^{v_n} \hat{a_n} b = T^{v_n} \hat{a_n} \otimes b = a_n\otimes b,$$
	implying by uniqueness of limits that $1\otimes ab = a\otimes b$ as $n\to \infty$, as desired.
\end{proof}

\begin{lemma}[$A_\infty$ is idempotent]
	\label{lemma: Ainfty is idempotent, f.g. case}
	
	The $A$-algebra $A_\infty$ is idempotent, by which we mean that the multiplication map of $A_\infty$ gives rise to an isomorphism $A_\infty\Dotimes_{(A, R)_\blacksquare} A_\infty \to A_\infty$ in $\derived{(A, R)_\blacksquare}$.
\end{lemma}
\begin{proof}
	By lemma \ref{lemma: formal Laurent series are idempotent} we know that multiplication $A_\infty\otimes_A A_\infty\to A_\infty$ is an $A$-module isomorphism. Now as $A_\infty$ is $A$-flat by lemma \ref{lemma: Ainfty is A-flat, f.g. case} the natural isomorphism $A_\infty\Dotimes_A A_\infty \to A_\infty\otimes_A A_\infty$ is a isomorphism. Recall that the derived completion $\derivedCompletion{-}{(A,R)_\blacksquare}$ is symmetric monoidal for $\Dotimes_{(A,R)_\blacksquare}$ by lemma \ref{lemma: completed scalar extension is monoidal}. But $A_\infty$ and hence $A_\infty \Dotimes_A A_\infty$ are $(A, R)_\blacksquare$-complete by lemma \ref{lemma: resolution for Ainfty, f.g. case} and thus in total
		$$A_\infty\Dotimes_{(A, R)_\blacksquare} A_\infty\cong A_\infty \Dotimes_A A_\infty \cong A_\infty\otimes_A A_\infty.$$
	Since the multiplication maps of all three tensor products are compatible and since the underived one is an isomorphism we obtain that
		$$A_\infty \Dotimes_{(A, R)_\blacksquare} A_\infty \to A_\infty$$
	is an isomorphism as well.
\end{proof}

An essentially purely formal consequence of the previous lemma \ref{lemma: Ainfty is idempotent, f.g. case} is given in the following lemma that characterizes $A_\infty$-modules inside of $\derived{(A, R)_\blacksquare}$. While Clausen and Scholze state this result, they omit the reasoning and provide no reference. The lemma follows from the fact that the \quote{Eilenberg--More-category of an \emph{idempotent} monad on a category $\category{C}$ is a reflective subcategory of $\category{C}$}. Alternatively this follows from work of Boyarchenko and Drinfeld in \cite{boyarchenko-drinfeld2009idempotents}, which rephrases this statement about monads in terms of the more familiar theory of monoidal categories.
\begin{lemma}[$A_\infty$-modules]
	\label{lemma: Ainfty-modules, f.g. case}
	
	The $A_\infty$-modules in $\derived{(A, R)_\blacksquare}$ form a full subcategory whose inclusion has a left adjoint given by the functor $A_\infty\Dotimes_{(A, R)_\blacksquare} -$. Furthermore, any $M \in \derived{(A, R)_\blacksquare}$ admits at most one $A_\infty$-module structure and such a structure exists if and only if the natural morphism $M\to A_\infty\Dotimes_{(A, R)_\blacksquare} M$ is an isomorphism.
\end{lemma}
\begin{proof}
	Denote by $\mu\colon A_\infty\Dotimes_{(A, R)_\blacksquare} A_\infty\to A_\infty$ the multiplication. In the language of \cite[Section 2]{boyarchenko-drinfeld2009idempotents}, the pair $e\ldef (A_\infty, \mu)$ is a \emph{unital algebra} in the monoidal category $\category{M}\ldef \derived{(A, R)_\blacksquare}$. Even more, since $\mu$ is an isomorphism, the pair is an \emph{idempotent} unital algebra. In particular the \emph{unit} $A\to A_\infty$ is an \emph{idempotent arrow} by \cite[Remark 2.9 (vi)]{boyarchenko-drinfeld2009idempotents}, making $(A_\infty, \mu)$ a \emph{closed} unital idempotent algebra. Denote by $\mod{e}$ the category of \emph{unital (left) $e$-modules} of this algebra.
	
	
	
	Consider the essential image $e\category{M}\ldef \{M \in \derived{(A, R)_\blacksquare}\setseparator A_\infty\Dotimes_{(A, R)_\blacksquare} M \cong M\}$ of the functor $M \mapsto A_\infty\Dotimes_{(A, R)_\blacksquare} M$ and denote by $L$ its (co)restriction $\category{M}\to e\category{M}$.
	
	By \cite[Lemma 2.32 (ii)]{boyarchenko-drinfeld2009idempotents} the forgetful functor $\mod{e} \to \category{M}$ is fully faithful and its essential image is $e\category{M}$, that is $\mod{e}\equiv e\category{M}$. A (pseudo)inverse $e\category{M} \to \category{M}$ is given by the obvious way of enriching $A_\infty\Dotimes_{(A, R)_\blacksquare} M$ with a unital $e=(A_\infty, \mu)$-module structure. Furthermore by \cite[Prop. 2.22 (a)]{boyarchenko-drinfeld2009idempotents} the functor $L\colon \category{M}\to e\category{M}$ is left adjoint to the fully faithful inclusion $e\category{M} \to \category{M}$, exhibiting $e\category{M}\equiv \mod{e}$ as a localization of $\category{M}$. Thus the functor $\mod{e} \to \derived{(A, R)_\blacksquare}$ is fully faithful with left adjoint given by $A_\infty\Dotimes_{(A, R)_\blacksquare} -$ as claimed. By \cite[Lemma 2.15]{boyarchenko-drinfeld2009idempotents} it follows that an $M$ admits an $(A_\infty, \mu)$-module structure if and only $M \to A_\infty\Dotimes_{(A, R)_\blacksquare} M$ is an isomorphism. As $e\category{M}\to \category{M}$ is fully faithful it is also clear that any $M\in\category{M}$ admits at most one unital $(A_\infty, \mu)$-module structure because if there are two modules $X, Y \in e\category{M}$ mapping both to $M \in \category{M}$ then the identity $\id_M$ lifts to an isomorphism $X \to Y$ of unital $(A_\infty, \mu)$-modules.
		
	It remains to see that the notion of $A_\infty$-module and unital $(A_\infty, \mu)$-module inside $\derived{(A, R)_\blacksquare}$ agree. For this observe that $A_\infty\Dotimes_{(A, R)_\blacksquare} M = \derivedCompletion{A_\infty\Dotimes_A M}{(A, R)_\blacksquare}$ since derived completion is (strong) symmetric monoidal by lemma \ref{lemma: completed scalar extension is monoidal}. Since further the completion is left adjoint to the inclusion into $A$-modules we find that
	$$\Hom_\derived{(A, R)_\blacksquare}(A_\infty\Dotimes_{(A, R)_\blacksquare} M, M) \cong \Hom_\derived{A}(A_\infty\Dotimes_A M, M).$$
	But since $A_\infty$ is $A$-flat by lemma \ref{lemma: Ainfty is A-flat, f.g. case} the natural morphism $A_\infty\Dotimes_A M \to A_\infty\otimes_A M$ is a quasi-isomorphism, yielding further that
	$$\Hom_\derived{(A, R)_\blacksquare}(A_\infty\Dotimes_{(A, R)_\blacksquare} M, M) \cong \Hom_\derived{A}(A_\infty \otimes_A M, M).$$
	Under this correspondence it is clear that inside of $\derived{(A, R)_\blacksquare}$ the notions of unital $(A_\infty, \mu)$-module and $A_\infty$-module agree, as one can translate the $A_\infty$-action on $M$ from one notion to the other.
	
%	
%	\begin{todo}
%		That any $A_\infty$-free module is complete is not clear.
%		
%		It would be much clearer to just show that any $A$-module only admits one honest/traditional $A_\infty$-module structure. Fullness is then clear as any A-linear map is also $A_\infty$-linear and faithfulness as scalar restriction is faithful already. The other statements then follow formally.
%	\end{todo}
%	By lemma \ref{lemma: resolution for Ainfty} the scalar restriction of $A_\infty$ and hence of any $A_\infty$-free module to $A$ is $(A, R)_\blacksquare$-complete. Thus the inclusion $A \subseteq A_\infty$ is a map of analytic rings $i\colon (A, R)_\blacksquare \to A_\infty$. By proposition \ref{proposition: scalar restriction and extension of complete modules} we obtain that the scalar restriction of any $A_\infty$-module is $(A, R)_\blacksquare$-complete and that we have an adjunction
%		$$\derived{A_\infty} \xleftrightarrows{\LeftD\completedScalarExtension{i}{A_\infty}}{\scalarRestriction{i}} \derived{(A, R)_\blacksquare}.$$
%	As $A_\infty$-completion is trivial, we obtain that $\LeftD\completedScalarExtension{i}{A_\infty} \cong \LeftD\scalarExtension{i}$ is simply the left derived functor of ordinary scalar extension. In particular $\scalarRestriction{i}$ is fully faithful if and only if the counit $\LeftD\scalarExtension{i}\scalarRestriction{i}\Rightarrow \id$ is an isomorphism, that is for any $M^\bullet\in\derived{A_\infty}$ the natural morphism $A_\infty\Dotimes_A M^\bullet \to M^\bullet$ is an isomorphism.
%	
%	Observe that since multiplication $A_\infty\otimes_A A_\infty \to A_\infty$ is an isomorphism by lemma \ref{lemma: Ainfty is idempotent}, so is the natural morphism $A_\infty\otimes_A \free{A_\infty}{S} \to \free{A_\infty}{S}$ for any extremally disconnected $S$. Let $M^\bullet \in \derived{A_\infty}$ be arbitrary. As $\scalarRestriction{i}$ and $\scalarExtension{i}$ are exact (by flatness of $A_\infty$), both $\scalarRestriction{i}$ and $\LeftD\scalarExtension{i}$ commute with limits, we can assume $M^\bullet$ to be bounded above, by writing it as a limit of its stupid truncations $\sTruncation{n}M^\bullet$. We can then choose a resolution $F^\bullet\to M^\bullet$ of $M^\bullet$ by $A_\infty$-free modules. Then all terms of $F^\bullet$ are $A$-flat as any $\free{A_\infty}{S} \cong A_\infty\otimes_A \free{A}{S}$ is a tensor product of flat modules by lemma \ref{lemma: Ainfty is A-flat}, hence flat. Then $F^\bullet$ is an $A$-flat resolution of $M^\bullet$ and we find that $A_\infty\Dotimes_A M^\bullet \to A_\infty \otimes_A F^\bullet$ is an isomorphism. But clearly $A_\infty\otimes_A F^\bullet \to F^\bullet$ is an isomorphism by the previous calculation. Thus $A_\infty\Dotimes_A M^\bullet\to M^\bullet$ is an isomorphism for all $M^\bullet \in \derived{A_\infty}$, showing the full faithfulness.
%	
%	
%	Suppose that an $A$-module $M^\bullet$ admits two $A_\infty$-multiplications $\cdot_1^\bullet$ and $\cdot_2^\bullet$. By full faithfulness $\id_{M^\bullet}$ provides an isomorphism between the resulting $A_\infty$-modules. Then for any $a\in A_\infty$, any $n\in \Z$ and any $m\in M^n$ we have
%		$$a\cdot_1^n m = \id_{M^n}(a\cdot_1^n m) = a\cdot_2^n \id_{M^n}(m) = a\cdot_2^n m$$
%	giving the desired $\cdot_1^\bullet = \cdot_2^\bullet$.
%	
%	\begin{todo}
%		This is ugly. Maybe also only need up to iso.
%		
%		"Exists if and only if this morphism is iso"
%	\end{todo}
\end{proof}

\begin{lemma}
	\label{lemma: Ainfty-modules admit no extensions, f.g. case}
	
	We have
		$$\eRHom_A(A_\infty, A) \cong 0$$
	in $\derived{\condensedAb}$. In particular if $C^\bullet \in \derived{\associated{A}}$ is a complex of $A_\blacksquare$-free modules and if $M$ is any $A_\infty$-module, then $\eRHom_A(M, C^\bullet)\cong 0$.
\end{lemma}
\begin{proof}
	Using the quasi-isomorphism of $A$-modules from lemma \ref{lemma: resolution for Ainfty, f.g. case}, we obtain that the complex $\eRHom_A(A_\infty, A)$ is quasi-isomorphic to the complex
	$$0\to \eRHom_R(\series{R}{U}, A) \xrightarrow{UT-1} \eRHom_R(\series{R}{U}, A)\to 0$$
	as $A\otimes_R \series{R}{U} \cong A\Dotimes_R \series{R}{U}$ by $R$-flatness of $A$ and since the $\condensedAb$-enriched adjunction of scalar extension and restriction is derivable by lemma \ref{lemma: enriched adjunctions are stable under derivation}. Clausen and Scholze now claim in \cite[Observation 8.8, p.55]{condenseddotpdf} that this complex is quasi-isomorphic to the complex
		$$0\to A[U^{-1}]/A \xrightarrow{UT-1} A[U^{-1}]/A\to 0,$$
	the author was however not able to verify this. We thus have to make this an \assumptionPreWord. Now this complex is acyclic since left multiplication by $UT-1$ is injective and the cokernel is $(A[U^{-1}]/A)/\ideal{UT-1}\cong A[T]/A\cong 0$ implying that $UT-1$ is surjective as well. So if we take the alternative representation as given we find that $\eRHom_R(A_\infty, A) \cong 0$ as claimed.
	
	Suppose that $C^\bullet = C[0]$ is concentrated in degree $0$ where $C = \bigoplus_{i\in I} \free{A_\blacksquare}{S_i}$ is $A_\blacksquare$-free. By Specker's theorem \ref{theorem: specker's theorem} there are sets $J_i$ for $i\in I$ such that $\free{A_\blacksquare}{S_i} \cong \prod_{J_i} A$. By lemma \ref{lemma: resolution for Ainfty, f.g. case} we find that $A_\infty$ is a derived compact object of $\derived{(A, R)_\blacksquare}$ so that by lemma \ref{lemma: solid derived compacts are enriched compact} it is also $\condensedAb$-enriched compact. As $\eRHom_A(A_\infty, -)$ additionally commutes with direct products as a consequence of remark \ref{remark: compatibility of enriched Homs with direct sums and products} we obtain that
		$$\eRHom_A(A_\infty, C^\bullet) \cong \bigoplus_{i\in I} \prod_{J_i} \eRHom_A(A_\infty, A)\cong 0$$
	as all these $A$-modules are $(A, R)_\blacksquare$-complete. Indeed $A_\infty$ is complete by lemma \ref{lemma: resolution for Ainfty, f.g. case} and the other cases follow since $A \cong A\otimes_R\free{R_\blacksquare}{\ast}$ is complete and since $\mod{(A, R)_\blacksquare}$ has all limits and colimits. By induction on the length and by utilizing stupid truncations it follows that any bounded complex $C^\bullet$ of $A_\blacksquare$-free modules also satisfies $\eRHom_A(A_\infty, C^\bullet) \cong 0$. Now if $C^\bullet\in\derived{A_\blacksquare}$ is bounded to the right then there is a left resolution $P^\bullet \to C^\bullet$ by $A_\blacksquare$-free modules which is also bound to the right. Writing $P^\bullet$ as a sequential colimit of its truncations $P_n^\bullet \ldef \sTruncation{\ge -n} P^\bullet$ we obtain that
		$$\eRHom_A(A_\infty, C^\bullet) \cong \eRHom_A(A_\infty, P^\bullet) \cong \colimit_{n\in \N} \eRHom_A(A_\infty, P_n^\bullet) \cong 0$$
	by lemma \ref{lemma: enriched compact objects preserve sequential homotopy colimits} as again, $A_\infty$ is enriched compact. Finally if $C^\bullet \in \derived{A_\blacksquare}$ is general, we can write it as the colimit $\colimit_{n\in \N} \sTruncation{\le n}C^\bullet$ of bound to the right complexes which implies
	$$\eRHom_A(A_\infty, C^\bullet)\cong \colimit_{n\in \N} \eRHom_A(A_\infty, \sTruncation{\le n} C^\bullet) \cong 0$$
	as $\eRHom_A(A_\infty, -)$ still preserves (homotopy) colimits by the same lemma.
	
	Now suppose that $M$ is an $A_\infty$-module. Then $M\cong M\Dotimes_{(A, R)_\blacksquare} A_\infty$ by lemma \ref{lemma: Ainfty-modules, f.g. case}. Thus for any $C^\bullet \in \derived{A_\blacksquare}$ we find that
	\begin{align*}
		\eRHom_A(M, C^\bullet) & \cong \eRHom_A(M\Dotimes_{(A, R)_\blacksquare} A_\infty, C^\bullet)\\
		&\cong \eRHom_A(\derivedCompletion{M\Dotimes_A A_\infty}{(A, R)_\blacksquare}, C^\bullet) \cong \eRHom_A(M\Dotimes_A A_\infty, C^\bullet).
	\end{align*}
	as derived completion is strong symmetric monoidal and as the $\condensedAb$-enriched adjunction between completion and inclusion into $A$-modules is derivable by \ref{lemma: enriched adjunctions are stable under derivation}. This then implies that
		$$\eRHom_A(M, C^\bullet) \cong \eRHom_A(M, \intRHom_A(A_\infty, C^\bullet))$$
	by lemma \ref{lemma: deriving a closed monoidal structure} as $\eRHom_A$ is the (complex of) condensed abelian group(s) underlying $\intRHom_A$. But as previously shown $\intRHom_A(A_\infty, C^\bullet) \cong 0$ as the underlying (complex of) condensed abelian group(s) of $\intRHom_A$ is $\eRHom_A$. Thus
		$$\eRHom_A(M,C^\bullet) \cong \eRHom_A(M, 0) \cong 0$$
	as well.
\end{proof}


\begin{lemma}
	\label{lemma: special cokernels are Ainfty-modules, f.g. case}
	
	For any set $I$, the cokernel of the injective map
	$$A\otimes_R \prod_I R \to \prod_I A$$
	is an $A_\infty$-module. In particular, by Specker's theorem \ref{theorem: specker's theorem} for any extremally disconnected $S$ the cokernel of the map $\free{(A, R)_\blacksquare}{S} \to \free{A_\blacksquare}{S}$ is an $A_\infty$-module.
\end{lemma}
\begin{proof}
	Recall that by lemma \ref{lemma: Ainfty-modules, f.g. case} any $A$-module admits at most one $A_\infty$-module structure and that an additive map between $A_\infty$-modules is $A_\infty$-linear if and only if it is $A$-linear.
	
	The short exact sequence
		$$0\to \prod_I R\to\prod_I \series{R}{T^{-1}} \to \prod_I T^{-1}\series{R}{T^{-1}}\to 0$$
	gives rise to the short exact sequence
		$$0\to A\otimes_R \prod_I R\to A\otimes_R \prod_I \series{R}{T^{-1}} \to A\otimes_R \prod_I T^{-1}\series{R}{T^{-1}} \to 0$$
	as $A=R[T]$ is $R$-flat. Observe that as $A$-modules, we have $A_\infty = \laurent{R}{T^{-1}} \cong \series{R}{T^{-1}}[T] \cong A \otimes_R \series{R}{T^{-1}}$ so that
		$$A_\infty \otimes_{\series{R}{T^{-1}}} \prod_I \series{R}{T^{-1}}\cong A\otimes_R \series{R}{T^{-1}}\otimes_{\series{R}{T^{-1}}} \prod_I \series{R}{T^{-1}}\cong A\otimes_R \prod_I \series{R}{T^{-1}},$$
	allowing us to replace the central term of the previous short exact sequence by an $A_\infty$-module.	Observe further that $A_\infty/A \cong T^{-1}\series R {T^{-1}}$. We obtain the following commutative diagram of condensed $A$-modules
	$$\begin{tikzcd}
		&0\ar[d] &0\ar[d] &0\ar[d]\\
		0\ar[r] &A\otimes_R \prod_I R \ar[r]\ar[d] &\prod_I A\ar[r]\ar[d] &C \ar[r]\ar[d] &0\\
		0\ar[r] &A_\infty\otimes_{\series{R}{T^{-1}}} \prod_I \series R {T^{-1}} \ar[r]\ar[d] &\prod_I A_\infty \ar[r]\ar[d] &C' \ar[r]\ar[d] &0\\
		0 \ar[r] &\prod_I T^{-1}\series R{T^{-1}} \ar[r]\ar[d] &\prod_I T^{-1}\series{R}{T^{-1}} \ar[r]\ar[d] &0 \ar[r]\ar[d] &0\\
		&0 &0 &0
	\end{tikzcd}$$
	with $C$ and $C'$ cokernels, such that all rows are exact. Then by the above argument $2$ out of the $3$ columns are exact and by the $9$ lemma so is the third, in particular $C\cong C'$ as $A$-modules. Now $C'$ is a cokernel of $A$-modules that admit a (then necessarily unique) $A_\infty$-module structure and the cokernel of this (necessarily $A_\infty$-linear) morphism is an $A_\infty$-module as well. Hence $C\cong C'$ is an $A_\infty$-module.
\end{proof}


We can now finally state the proof of the analogue to Hilbert's Basis Theorem.
\begin{proof}[of theorem \ref{theorem: a kind of hilbert's basis theorem}]
	Let $C^\bullet$ be a complex of $A_\blacksquare$-free modules concentrated in non-negative degrees. Let $S$ be extremally disconnected. By Specker's theorem \ref{theorem: specker's theorem} each term of $C^\bullet$ is a direct sum of products of copies of $A$. As $A\cong A\otimes_R \free {R_\blacksquare}{\ast}$ is $(A, R)_\blacksquare$-complete, so is each term of $C^\bullet$ and hence $C^\bullet\in\derived{(A,R)_\blacksquare}$. Thus
		$$\eRHom_A(\free{(A, R)_\blacksquare}{S}, C^\bullet) \isorightarrow \eRHom_A(\free A S, C^\bullet).$$
	By lemma \ref{lemma: special cokernels are Ainfty-modules, f.g. case} the cokernel $M$ of the natural morphism $\free{(A,R)_\blacksquare}{S}\to\free{A_\blacksquare}{S}$ is an $A_\infty$-module. We obtain the distinguished triangle
		$$\eRHom_A(M, C^\bullet)\to\eRHom_A(\free{A_\blacksquare}{S}, C^\bullet)\to\eRHom_A(\free{(A, R)_\blacksquare}{S}, C^\bullet)\to (\cdots)[1]$$
	for which $\eRHom_A(M, C^\bullet) \cong 0$ by lemma \ref{lemma: Ainfty-modules admit no extensions, f.g. case}. Hence
		$$\eRHom_A(\free{A_\blacksquare}{S}, C^\bullet)\to\eRHom_A(\free{(A, R)_\blacksquare}{S}, C^\bullet)$$
	is an isomorphism. Overall we thus obtain that the natural morphism
	$$\eRHom_A(\free{A_\blacksquare}{S}, C^\bullet)\isorightarrow \eRHom_A(\free{(A, R)_\blacksquare}{S}, C^\bullet) \isorightarrow \eRHom_A(\free A S, C^\bullet)$$
	is an isomorphism. Thus any such complex $C^\bullet$ of $A_\blacksquare$-free modules is derived $A_\blacksquare$-complete and hence $A_\blacksquare$ is analytic.
\end{proof}
